\subsection{真偏映射}

在本节中，X 与 Y 是局部紧拓扑空间。我们用 \(X'\)（相应地，\(Y'\)）表示在 X（相应地，Y）上添加一个无穷远点 \(\omega_{X}\)（相应地，\(\omega_{Y}\)）所得到的紧空间 (TG, I, p.~67--68)。我们把 \(X'\) 与 \(Y'\) 分别等同于 \(\mathsf{X}'(\mathcal{C}_{0}(\mathsf{X}))\) 与 \(\mathsf{X}'(\mathcal{C}_{0}(\mathsf{Y}))\) (I, p.~32, 推论 2)。

定义 1. --- 从 X 到 Y 的真偏映射是指 X 与 Y 之间的一个对应 \(f = (\Gamma, \mathrm{X}, \mathrm{Y})\) (E, II, p.~10, 定义 2)，它满足：

\begin{enumerate}
\def\labelenumi{(\roman{enumi})}
\item
  图像 \(\Gamma\) 是函数性的；
\item
  f 的定义域是 X 的一个开集 U；
\item
  从 U 到 Y 的映射 \(x\mapsto f(x)\) 是真映射。
\end{enumerate}

X 上的恒等映射是从 X 到 X 的真偏映射。设 Z 是一个局部紧拓扑空间，并设 \(f\)（相应地，\(g\)）是从 X 到 Y（相应地，从 Y 到 Z）的真偏映射。则复合对应 \(g \circ f\) (E, II, p.~11, 定义 6) 是从 X 到 Z 的真偏映射 (TG, I, p.~72, 命题 3 及 p.~73, 命题 5)。

引理 1. --- 对每个以 U 为定义域的、从 X 到 Y 的真偏映射 f，用 \(\widetilde{f}\) 表示由 \(\widetilde{f}(x) = f(x)\)（若 \(x \in \mathrm{U}\)）与 \(\widetilde{f}(x) = \omega_{\mathrm{Y}}\)（若 \(x \notin \mathrm{U}\)）定义的从 X' 到 Y' 的映射；它是连续的。

映射 \(f \mapsto \widetilde{f}\) 是从 X 到 Y 的真偏映射 \(f\) 所成的集合与从 X' 到 Y' 的连续映射 \(g\)（满足 \(g(\omega_{\mathrm{X}}) = \omega_{\mathrm{Y}}\)）所成的集合之间的一个双射。

设 \(f\) 是从 X 到 Y 的一个真偏映射，并设 U 是它的定义域。我们来证明映射 \(\widetilde{f}\) 是连续的。它在 U 的每一点都连续，因为 U 在 \(\mathrm{X}'\) 中是开的。我们来证明它在每一点 \(x\) ∈ \(\mathrm{X}' - \mathrm{U}\) 处也连续；这时有 \(\widetilde{f}(x) = \omega_{\mathrm{Y}}\)。设 V 是 \(\omega_{\mathrm{Y}}\) 在 \(\mathrm{Y}'\) 中的一个开邻域；我们来证明 \(\widetilde{f}^{-1}(\mathrm{V})\) 是 \(x\) 的一个邻域。根据拓扑空间 \(\mathrm{Y}'\) 的定义，可以假定 V 形如 \(Y^{\prime}-K\)，其中 K 是 Y 的一个紧子集。由于 f 定义了从 U 到 Y 的一个真映射，集合 \(f^{-1}(K)\) 在 U 中是紧的 (TG, I, p.~77, 命题 6)，因而在 \(X^{\prime}\) 中是紧的。它特别是 \(X^{\prime}\) 的一个闭子集，于是 \(\tilde{f}^{-1}(V)=X^{\prime}-f^{-1}(K)\) 是 \(X^{\prime}\) 的一个开子集，从而是 x 的一个邻域。

反之，设 \(g \colon X' \to Y'\) 是满足 \(g(\omega_{\mathrm{X}}) = \omega_{\mathrm{Y}}\) 的连续映射，并设 \(\Gamma_g \subset X' \times Y'\) 是它的图像。集合 \(U = X - \overline{g}^{-1}(\omega_{\mathrm{Y}})\) 在 \(X\) 中是开的。对应 \(f = (\Gamma_g \cap (U \times Y), X, Y)\) 是从 \(X\) 到 \(Y\) 的一个真偏映射，它满足 \(\widetilde{f} = g\)，并且是唯一的一个。

我们将把从 X 到 Y 的真偏映射与从 \(X'\) 到 \(Y'\) 中的、把 \(\omega_{X}\) 映到 \(\omega_{Y}\) 的连续映射视为同一。特别地，从 X 到 Y 的真映射就是以 X 为定义域的真偏映射；它们与从 \(X'\) 到 \(Y'\) 中的、满足 \(f^{-1}(\omega_{\mathrm{Y}}) = \{\omega_{\mathrm{X}}\}\) 的连续映射 f 视为同一。若 X 是紧的，这些映射无非就是从 X 到 Y 的连续映射。

设 A 是一个复交换 Banach 代数。回顾 \(X'(A)\) 与在 \(X(A)\) 上添加一个无穷远点所得到的紧空间视为同一 (I, p.~29, 推论)。

命题 2. --- 设 A 与 B 是复交换 Banach 代数。对每个代数同态 \(\pi\colon\) A \(\to\) B，映射 \(\mathsf{X}'(\pi)\) 是从 \(\mathsf{X}(\mathsf{B})\) 到 \(\mathsf{X}(\mathsf{A})\) 的真偏映射。

事实上，\(\mathsf{X}'(\pi)\) 是从 \(\mathsf{X}'(\mathrm{B})\) 到 \(\mathsf{X}'(\mathrm{A})\) 的连续映射 (I, p.~10)。\(\mathsf{X}'(\mathrm{B})\)（相应地，\(\mathsf{X}'(\mathrm{A})\)）的无穷远点是零特征标，并且我们有 \(\mathsf{X}'(\pi)(0)=0\)。

命题 3. --- a) 对每个从 X 到 Y 的真偏映射 \(\varphi\)，映射 \(f \mapsto f \circ \varphi\)（从 \(\mathcal{C}(\mathrm{Y}')\) 到 \(\mathcal{C}(\mathrm{X}')\)）诱导出一个代数态射 \(\varphi^*\)（从 \(\mathcal{C}_0(\mathrm{Y})\) 到 \(\mathcal{C}_0(\mathrm{X})\)）；

\begin{enumerate}
\def\labelenumi{\alph{enumi})}
\setcounter{enumi}{1}
\tightlist
\item
  映射 \(\varphi \mapsto \varphi^{*}\) 是从 X 到 Y 的真偏映射所成的集合到由 \(\mathcal{C}_0(\mathrm{Y})\) 到 \(\mathcal{C}_0(\mathrm{X})\) 的代数态射所成的集合上的一个双射。其逆双射是映射 \(\pi \mapsto \mathsf{X}'(\pi)\)。
\end{enumerate}

我们来证明 a)。设 \(\varphi\) 是从 X 到 Y 的真偏映射，把它与满足 \(\varphi(\omega_{\mathrm{X}}) = \omega_{\mathrm{Y}}\) 的、从 X' 到 Y' 的连续映射视为同一。对 \(f \in \mathcal{C}_0(\mathrm{Y})\)，有 \((f \circ \varphi)(\omega_{\mathrm{X}}) = f(\omega_{\mathrm{Y}}) = 0\)，因而映射 \(\varphi^*\) 是良定义的。它是一个代数态射。

我们来证明 \(\mathsf{X}'(\varphi^{*})\) 与 \(\varphi\) 视为同一。设 \(x \in X\)。对每个函数 \(f \in \mathcal{C}_{0}(\mathrm{Y})\)，特征标 \(\mathsf{X}'(\varphi^{*})(\mathrm{ev}_{x})\) 把复数

\[
\left(\operatorname{ev} _ {x} \circ \varphi^ {*}\right) (f) = \operatorname{ev} _ {x} (f \circ \varphi) = f (\varphi (x)),
\]

与 f 相联系，因而 \(\mathsf{X}'(\varphi^{*})(\mathrm{ev}_{x}) = \mathrm{ev}_{\varphi(x)}\)。这就证明了该断言。

反之，设 \(\pi\colon\mathcal{C}_{0}(\mathrm{Y})\to\mathcal{C}_{0}(\mathrm{X})\) 是一个代数态射。我们来证明 \(\mathsf{X}^{\prime}(\pi)^{*}=\pi\)。设 \(f\in\mathcal{C}_{0}(\mathrm{Y})\)，并记 \(g=\mathsf{X}^{\prime}(\pi)^{*}(f)\in\mathcal{C}_{0}(\mathrm{X})\)。对每个 \(x\in\mathsf{X}\)，有

\[
g (x) = \left(f \circ \mathrm{X} ^ {\prime} (\pi)\right) (x) = \operatorname{ev} _ {\mathrm{X} ^ {\prime} (\pi) (x)} (f) = \left(\operatorname{ev} _ {x} \circ \pi\right) (f) = \pi (f) (x),
\]

因为 \(\mathsf{X}^{\prime}(\pi)\) 满足 \(\mathrm{ev}_{\mathsf{X}^{\prime}(\pi)(x)} = \mathrm{ev}_{x} \circ \pi\)。于是我们有 \(g = \pi(f)\)，由此得出 \(\mathsf{X}^{\prime}(\pi)^{*} = \pi\)。

以完全类似的方式，我们有：

命题 4. --- 假定 X 与 Y 是紧的。把空间 X（相应地，空间 Y）与 X( \(\mathcal{C}(X)\) )（相应地，X( \(\mathcal{C}(Y)\) )）视为同一 (I, p.~33, 推论 3)。

\begin{enumerate}
\def\labelenumi{\alph{enumi})}
\item
  对每个连续映射 \(\varphi\colon\mathrm{X}\to\mathrm{Y}\)，映射 \(\varphi^{*}\colon f\mapsto f\circ\varphi\) 是从 \(\mathcal{C}(\mathrm{Y})\) 到 \(\mathcal{C}(\mathrm{X})\) 的代数态射；
\item
  映射 \(\varphi \mapsto \varphi^{*}\) 与 \(\pi \mapsto X(\pi)\) 是从 \(X\) 到 \(Y\) 的连续映射所成的集合与从 \(\mathcal{C}(Y)\) 到 \(\mathcal{C}(X)\) 的代数态射所成的集合之间互逆的双射。
\end{enumerate}

注. --- *用范畴论的语言，前面的结果可解释如下。设 G 是这样一个范畴：其对象为局部紧拓扑空间，其态射为真偏映射。函子 X ↦ \(\mathcal{C}_0(X)\) 是从范畴 G 到复交换 Banach 代数范畴的一个反变的全忠实函子。此外，A ↦ X(A) 是从复交换 Banach 代数范畴到范畴 G 的一个反变函子。若对每个局部紧拓扑空间 X 联系同胚

\[
\operatorname{ev} \colon X \to X (\mathscr {C} _ {0} (X)),
\]

则得到从范畴 G 的恒等函子到复合函子 \(X \mapsto X(\mathcal{C}_{0}(X))\) 的一个同构。

复合函子 \(A \mapsto \mathcal{C}_{0}(X(A))\) 并不同构于复交换 Banach 代数范畴的恒等函子（参见 I, p.~36, 例 2 及 I, p.~155, 习题 2）。不过，对于交换星代数（C*-代数），我们将看到一个这种类型的命题 (I, p.~107, 第 5 小节）。*

\subsection{Gelfand 变换}

设 A 是一个交换 Banach 代数。回顾：对每个 \(x \in A\)，我们用 \(\mathcal{G}_{\mathrm{A}}(x)\) 或 \(\mathcal{G}(x)\) 表示函数 \(\chi \mapsto \chi(x)\)（定义在 \(\mathsf{X}(\mathsf{A})\) 上）；\(\mathcal{G}(x)\) 称为 x 的 Gelfand 变换，而映射 \(x \mapsto \mathcal{G}(x)\) 称为 Gelfand 变换（参见 I, p.~7, 定义 5）。于是按定义有：

\[
\mathcal {G} (x) (\chi) = \chi (x).
\]

例. --- 1) 设 X 是一个局部紧拓扑空间，并考虑交换 Banach 代数 \(\mathcal{C}_{0}(\mathrm{X})\) (I, p.~17, 例 3 及 I, p.~31, 第 2 小节)。由 I, p.~32, 推论 1，特征标空间 \(\mathsf{X}(\mathcal{C}_{0}(\mathbf{X}))\) 借助于把元素 \(x \in \mathbf{X}\) 联系到特征标 \(f \mapsto f(x)\)（\(\mathcal{C}_{0}(\mathbf{X})\) 的特征标）的映射而与 X 视为同一，于是 \(\mathcal{C}_{0}(\mathbf{X})\) 的 Gelfand 变换与恒等映射视为同一。

\begin{enumerate}
\def\labelenumi{\arabic{enumi})}
\setcounter{enumi}{1}
\tightlist
\item
  设 \(n \geqslant 0\) 是一个整数。设 \(A_n\) 是由函数 \(f: [0,1] \to K\)（在 \([0,1]\) 上具有直到 \(n\) 阶的连续导数）所成的代数。赋以范数
\end{enumerate}

\[
\| f \| = \sum_ {k = 0} ^ {n} \frac {1}{k !} \sup _ {0 \leqslant t \leqslant 1} | f ^ {(k)} (t) |,
\]

它是一个 Banach 代数 (I, p.~18, 例 4)。对每个 n，特征标空间 \(\mathsf{X}(\mathsf{A}_{n})\) 与 {[}0,1{]} 视为同一，而 G 与 \(A_{n}\) 到 \(\mathcal{C}([0,1])\) 中的包含映射视为同一（参见 I, p.~144, 例 1）。

\begin{enumerate}
\def\labelenumi{\arabic{enumi})}
\setcounter{enumi}{2}
\item
  设 \(\Delta\) 是由复数 \(z\)（满足 \(|z| \leqslant 1\)）组成的圆盘，并设 A 是由 \(\Delta\) 上连续且在 \(\Delta\) 的内部解析的函数所成的复 Banach 代数，赋以范数 \(\|f\| = \sup_{z \in \Delta} |f(z)|\) (I, p.~20, 例 9)。则 \(\mathsf{X}(\mathsf{A})\) 与 \(\Delta\) 视为同一，而 \(\mathcal{G}\) 与 A 到 \(\mathcal{C}(\Delta)\) 中的包含映射视为同一（参见 I, p.~193, 习题 6）。
\item
  考虑绝对收敛 Fourier 级数所成的复 Banach 代数 A (I, p.~19, 例 8)。对单位圆 U 的每个元素 \(u\)，映射 \(f \mapsto f(u)\) 是 A 的一个特征标 \(ev_{u}\)。若 \(f_{0} \in A\) 是 U 的恒等映射，则有 \(ev_{u}(f_{0}) = u\)，因而映射 ev: \(u \mapsto ev_{u}\)（从 U 到 X(A)）是单射；它是连续的。设 \(\chi \in \mathsf{X}(\mathsf{A})\)。我们有 \(\|f_{0}\| = \|f_{0}^{-1}\| = 1\)，因而 \(|\chi(f_{0})| \leqslant 1\) 且 \(|\chi(f_{0})^{-1}| \leqslant 1\)。这表明 \(\chi(f_{0}) \in \mathbf{U}\)，并且存在 \(u \in U\) 使得 \(\chi(f_{0}) = \mathrm{ev}_{u}(f_{0})\)。由于 \(\{f_{0}, f_{0}^{-1}\}\) 拓扑生成幺代数 A，有 \(\chi = ev_{u}\)。于是，映射 ev 是从 U 到 X(A) 上的一个同胚，借此把这两个空间视为同一。A 的 Gelfand 变换于是与 A 到 \(\mathcal{C}(\mathbf{U})\) 中的包含映射视为同一。
\end{enumerate}

由于 A 同构于 Banach 代数 \(\mathrm{L}^{1}(\mathbf{Z})\) (I, p.~19, 例 8)，空间 \(\mathsf{X}(\mathrm{L}^{1}(\mathbf{Z}))\) 与 U 视为同一，并且对每个元素 \((c_{n})\in\mathrm{L}^{1}(\mathbf{Z})\)，Gelfand 变换 \(\mathcal{G}_{\mathrm{L}^{1}(\mathbf{Z})}((c_{n}))\) 与 U 上的函数 \(u\mapsto\sum_{n\in\mathbf{Z}}c_{n}u^{n}\) 视为同一。

\begin{enumerate}
\def\labelenumi{\arabic{enumi})}
\setcounter{enumi}{4}
\tightlist
\item
  设 \(\Delta\) 是由复数 \(z\)（满足 \(|z| \leqslant 1\)）组成的单位圆盘。它在 \(\mathbf{C}\) 中的边界是 \(\mathbf{U}\)。设 A 是由复函数 \(f\)（\(\mathbf{U}\) 上的函数，存在连续函数 \(\widetilde{f} \in \mathcal{C}(\Delta)\) 延拓 \(f\) 且在 \(\mathring{\Delta}\) 中解析）所成的 Banach 代数，赋以范数 \(\|f\| = \sup_{z \in \mathbf{U}} |f(z)|\)。于是由最大模原理 (VAR, R1, p.~30, 3.3.7)，有 \(\|f\| = \sup_{z \in \Delta} |\widetilde{f}(z)|\)，因而 A 与 I, p.~20, 例 9 中的代数重合。集合 \(\mathsf{X}(\mathsf{A})\) 与 \(\Delta\) 视为同一，且若 \(f \in \mathsf{A}\)，则映射 \(\mathcal{G}(f)\) 与 \(f\) 到 \(\Delta\) 上的、在 \(\mathring{\Delta}\) 中解析的连续延拓视为同一（参见 I, p.~193, 习题 6）。
\end{enumerate}

命题 5. --- 设 A 是一个交换 Banach 代数。对每个 \(x \in \mathbf{A}\)，函数 \(\mathcal{G}(x)\) 属于交换 Banach 代数 \(\mathcal{C}_0(\mathsf{X}(\mathbf{A}))\)（由 \(\mathsf{X}(\mathbf{A})\) 上在无穷远处为 0 的连续函数所成）。

按定义（参见 I, p.~9, 第 7 小节），函数 \(\mathcal{G}_{\mathrm{A}}^{\prime}(x)\colon \chi \mapsto \chi (x)\) 在 \(\mathsf{X}'(\mathsf{A})\) 上连续且在 0 处取值 0。由于 \(\mathsf{X}'(\mathsf{A})\) 依 I, p.~29, 推论 1 与 \(\mathsf{X}(\mathsf{A})\) 的 Alexandroff 紧化视为同一，命题得证。

命题 6. --- 设 A 是一个交换 Banach 代数，并设 \(x \in A\)。

\begin{enumerate}
\def\labelenumi{\alph{enumi})}
\item
  \(\mathcal{G}(x)\) 的值集与 \(\{0\}\) 的并等于 \(\mathrm{Sp}_{\mathrm{A}}^{\prime}(x)\)；
\item
  若 A 有单位元，则 \(\mathcal{G}(x)\) 的值集就是 \(\mathrm{Sp}_{\mathrm{A}}(x)\)。特别地，要使 \(x\) 可逆，必须且只需 \(\mathcal{G}(x)\) 不取值 0。
\end{enumerate}

假定 A 有单位元。我们知道，对每个 \(\chi \in \mathsf{X}(\mathsf{A})\)，有 \(\chi(x) \in \mathrm{Sp}_{\mathsf{A}}(x)\)。反之，设 \(\lambda \in \mathrm{Sp}_{\mathsf{A}}(x)\)。则 \(x - \lambda\) 不可逆，因而属于 A 的某个极大理想。于是存在 \(\chi \in \mathsf{X}(\mathsf{A})\) 使得 \(\chi(x - \lambda) = 0\) (I, p.~30, 定理 2)，由此得 \(b\)）。

过渡到一般情形。设 \(\widetilde{\mathsf{A}}\) 是由 A 添加单位元所得到的 Banach 代数；它是交换的。集合 \(\mathrm{Sp}_{\widetilde{\mathsf{A}}}^{\prime}(x)\) 等于 \(\mathrm{Sp}_{\widetilde{\mathsf{A}}}^{\sim}(x)\)，即 \(\mathcal{G}_{\widetilde{\mathsf{A}}}^{\sim}(x)\) 在 \(\mathsf{X}(\widetilde{\mathsf{A}}) = \mathsf{X}'(\mathsf{A})\) 上的值集。由此得 \(a\)）。

例. --- 考虑绝对收敛 Fourier 级数所成的 Banach 代数 A (例 4)。命题 6, b) 蕴涵：若 \(\varphi\) 是单位圆 U 上一个具有绝对收敛 Fourier 级数的函数，且 \(\varphi\) 不取值 0，则函数 \(1/\varphi\) 也具有绝对收敛 Fourier 级数（“Wiener 定理”）。

命题 7. --- 设 A 是一个交换 Banach 代数。

\begin{enumerate}
\def\labelenumi{\alph{enumi})}
\item
  Gelfand 变换 \(\mathcal{G}\) 定义一个从 A 到 \(\mathcal{C}_0(\mathsf{X}(\mathsf{A}))\) 的态射，使得 \(\| \mathcal{G}(x)\| = \varrho (x)\leqslant \| x\|\) 对一切 \(x\in \mathrm{A}\) 成立；
\item
  要使 Gelfand 变换 \(\mathcal{G}\) 是等距的，必须且只需 \(\|x^{2}\|=\|x\|^{2}\) 对一切 \(x\in\mathrm{A}\) 成立。
\end{enumerate}

依 I, p.~9, 第 7 小节及 I, p.~37, 命题 5，映射 \(\mathcal{G}\) 是从 A 到 \(\mathcal{C}_0(\mathsf{X}(\mathsf{A}))\) 的态射；依命题 6 及 I, p.~28, 推论 5，它满足 \(\| \mathcal{G}(x)\| = \varrho (x)\)。断言 \(b)\) 由 \(a)\) 与 I, p.~21, 注 1 得出。

推论. --- 设 A 是一个 Banach 代数，并设 x 与 y 是 A 中可交换的元素。

\begin{enumerate}
\def\labelenumi{\alph{enumi})}
\item
  我们有 \(\varrho(xy) \leqslant \varrho(x)\varrho(y)\) 及 \(\varrho(x + y) \leqslant \varrho(x) + \varrho(y)\)；
\item
  若 y 是拟幂零元，则 \(\mathrm{Sp}_{\mathrm{A}}^{\prime}(x)=\mathrm{Sp}_{\mathrm{A}}^{\prime}(x+y)\)；此外若 A 是幺的，则 \(\mathrm{Sp}_{\mathrm{A}}(x)=\mathrm{Sp}_{\mathrm{A}}(x+y)\)。
\end{enumerate}

考虑由 A 添加单位元所得到的 Banach 代数，首先化归到代数 A 是幺的情形。然后，考虑由 x 与 y 生成的 A 的闭全子代数，就化归到 A 是交换且幺的情形。这时断言 (a) 是命题 7(a) 的推论，而断言 (b) 由 I, p.~37, 命题 6 及 I, p.~28, 推论 5 得出。

命题 8. --- 设 A 是一个交换 Banach 代数。下列四个集合相等：

\begin{enumerate}
\def\labelenumi{(\roman{enumi})}
\item
  Gelfand 变换的核；
\item
  A 中满足 \(\mathrm{Sp}_{\mathrm{A}}^{\prime}(x)=\{0\}\) 的元素 x 所成的集合；
\item
  A 的拟幂零元所成的集合；
\item
  A 的根。
\end{enumerate}

我们分别用 \(N_{1}\)、\(N_{2}\)、\(N_{3}\)、\(N_{4}\) 表示这些集合。我们有 \(N_{1} = N_{2}\)（命题 6, a)），以及 \(N_{2} = N_{3}\) (I, p.~28, 推论 5)。按定义，集合 \(N_{4}\) 是 A 的正则极大理想的交；因而它是 A 的诸特征标的核的交 (I, p.~30, 定理 2)，后者等于 \(N_{1}\)。

注. --- 1) 一般说来，Gelfand 变换的像既不在 \(\mathcal{C}_{0}(\mathsf{X}(\mathsf{A}))\) 中闭，也不在 \(\mathcal{C}_{0}(\mathsf{X}(\mathsf{A}))\) 中稠密 (I, p.~193, 习题 7)。

\begin{enumerate}
\def\labelenumi{\arabic{enumi})}
\setcounter{enumi}{1}
\item
  Gelfand 变换的像分离 \(\mathsf{X}(\mathsf{A})\) 的点，因为若 \(\chi_1 \neq \chi_2\) 是 \(\mathsf{X}(\mathsf{A})\) 的元素，则存在 \(x \in \mathsf{A}\) 使得 \(\chi_1(x) \neq \chi_2(x)\)。
\item
  若 \(\chi \in \mathsf{X}(\mathsf{A})\)，则存在 Gelfand 变换的像中的一个元素，它在 \(\chi\) 处不取值 0。
\item
  若 A 有单位元，则 Gelfand 变换的像是 \(\mathsf{X}(\mathsf{A})\) 上连续函数代数的一个全子代数（命题 6, b)）。
\end{enumerate}

引理 2. --- 设 A 是一个交换 Banach 代数。设 M 是 X(A) 的一个子集。则 M 关于 Jacobson 拓扑是闭的，当且仅当对每个特征标 \(\chi \in \mathsf{X}(\mathsf{A}) = \mathsf{M}\)，存在 A 的一个元素 \(x\)，使得 \(\mathcal{G}(x)\) 在 M 上为零而在 \(\chi\) 处非零。

设 \(\Upsilon(\mathrm{M})\) 是 M 的诸元素的核的交。集合 M 关于 Jacobson 拓扑是闭的，当且仅当 \(\mathrm{M} = \mathrm{V}(\Upsilon(\mathrm{M}))\)（参见 I, p.~13 及 I, p.~30）。这个条件等价于说：M 的元素 \(\chi\) 恰好是在 \(\Upsilon(\mathrm{M})\) 上为零的那些特征标。因此，M 是闭的当且仅当对每个特征标 \(\chi \notin \mathrm{M}\)，存在 \(x \in \Upsilon(\mathrm{M})\) 使得 \(\chi(x) \neq 0\)。这可转写为 \(\mathcal{G}(x)(\chi) \neq 0\) 且 \(\mathcal{G}(x)|_{\mathrm{M}} = 0\)。

\subsection{交换 Banach 代数的态射}

命题 9. --- 设 A 是一个 Banach 代数，并设 B 是一个交换半单 Banach 代数。从 A 的基础代数到 B 的基础代数的每个态射都是连续的。

设 \(h \colon \mathrm{A} \to \mathrm{B}\) 是一个代数同态，并设 \((a, b) \in \mathrm{A} \times \mathrm{B}\) 是附着于图像 \(\Gamma\)（\(h\) 的图像）的一个点。设 \(\chi \in \mathsf{X}'(\mathrm{B})\)。函数 \(x \mapsto \chi(h(x))\)（从 \(\mathrm{A}\) 到 \(\mathbf{C}\)）是一个代数同态，因而是连续的 (I, p.~29, 定理 1)。于是从 \(\mathrm{A} \times \mathrm{B}\) 到 \(\mathbf{C}\) 中的、由 \((x, y) \mapsto \chi(h(x)) - \chi(y)\) 给出的映射是连续的，它在 \(\Gamma\) 上为零，因而在 \((a, b)\) 处为零。这样，我们有 \(\chi(h(a)) = \chi(b)\) 对一切 \(\chi \in \mathsf{X}'(\mathrm{B})\) 成立。由于 B 是半单的，我们有 \(h(a) = b\)。因而 \(h\) 的图像是闭的，从而 \(h\) 是连续的 (EVT, I, p.~19, 推论 5)。

推论. --- 在一个无根的复交换代数上，定义 Banach 代数结构的两个范数是等价的。

只需对该代数的恒等映射应用命题 9。

设 A 与 B 是交换 Banach 代数。依 I, p.~9, 第 7 小节，若 \(h\colon A\to B\) 是一个满射态射，则 \(\mathsf{X}'(h)\) 是从 \(\mathsf{X}'(\mathsf{B})\) 到 \(\mathsf{X}'(\mathsf{A})\) 的一个闭子空间上的同胚，且它把 0 映到 0。（在 \(h\) 是一个子代数到 \(\mathcal{C}(\mathsf{X})\) 中的单射的情形，在弱得多的假设下映射 \(\mathsf{X}'(h)\) 就是单射，参见 I, p.~142, 命题 1, d)。）

现在设 \(h \colon \mathrm{A} \to \mathrm{B}\) 是一个单射态射。一般说来，\(\mathsf{X}'(h)\) 不是满射的，但下面的命题给出了使之成立的一个必要条件。

命题 10. --- 设 A 与 B 是幺交换 Banach 代数，\(h\colon A\to B\) 是一个保单位元的代数同态，不必连续。若 \(\mathsf{X}(h)\) 是满射的，则 \(h(\mathrm{A})\) 是 B 的一个全子代数。

设 \(x \in A\) 使得 \(h(x)\) 在 B 中可逆。对每个 \(\chi \in \mathsf{X}(\mathsf{A})\)，存在 \(\xi \in \mathsf{X}(\mathsf{B})\) 使得 \(\chi = \mathsf{X}(h)(\xi)\)，因而 \(\chi(x) = \xi(h(x)) \neq 0\)。于是 I, p.~37, 命题 6 表明 x 在 A 中可逆，从而 \(h(x)\) 在 \(h(\mathsf{A})\) 中可逆。

该命题中的必要条件不是充分的，即使 h 是等距的也是如此 (I, p.~168, 习题 14)。不过我们有如下的结果：

命题 11. --- 设 A 与 B 是交换幺 Banach 代数，a 是 A 的一个元素，而 h: A → B 是一个单的保单位元态射（不必连续）。假定由 a 生成的 A 的闭全子代数等于 A。下列条件等价：

\begin{enumerate}
\def\labelenumi{(\roman{enumi})}
\item
  \(\mathsf{X}(h)\) 是满射的；
\item
  \(h(\mathrm{A})\) 是 B 的一个全子代数；
\item
  \(\mathrm{Sp}_{\mathrm{A}}(a)=\mathrm{Sp}_{\mathrm{B}}(h(a)).\)
\item
  (i) \(\Longrightarrow\) (ii) 由命题 10 得出。
\item
  (ii) \(\Longrightarrow\) (iii) 由公式 \(\mathrm{Sp}_{\mathrm{A}}(a) = \mathrm{Sp}_{h(\mathrm{A})}(h(a)) = \mathrm{Sp}_{\mathrm{B}}(h(a))\) 得出，该公式有效是因为 \(h(\mathrm{A})\) 是 B 的一个全子代数。
\item
  (iii) \(\Longrightarrow\) (i) 依 I, p.~6, 第 6 小节的公式 (4)，我们有交换图表
\end{enumerate}

\[
\begin{array}{c} \mathsf {X} (\mathrm{B}) \xrightarrow {\mathsf {X} (h)} \mathsf {X} (\mathrm{A}) \\ \Big \downarrow^ {\mathcal {G} _ {\mathrm{B}} (h (a))} \qquad \qquad \qquad \Big \downarrow^ {\mathcal {G} _ {\mathrm{A}} (a)} \\ \operatorname{Sp} _ {\mathrm{B}} (h (a)) \xrightarrow {i} \operatorname{Sp} _ {\mathrm{A}} (a) \end{array}
\]

其中竖直箭头表示满射 (I, p.~37, 命题 6)，而 \(i\) 是典范包含。假设 (iii) 意味着 \(i\) 是双射。此外，满射 \(\mathcal{G}_{\mathrm{A}}(a)\colon \mathsf{X}(\mathrm{A})\to \mathrm{Sp}_{\mathrm{A}}(a)\) 是双射：事实上，对 A 的任意特征标 \(\chi_{1}\) 与 \(\chi_{2}\)，集合 \(\mathrm{A}_{\chi_1,\chi_2} = \{x\in \mathrm{A}\mid \chi_1(x) = \chi_2(x)\}\) 是 A 的一个闭全子代数。于是依假设，我们有 \(\mathrm{A}_{\chi_1,\chi_2} = \mathrm{A}\)（若 \(a\in \mathrm{A}_{\chi_1,\chi_2}\)），也就是说，\(\chi_{1} = \chi_{2}\)（若 \(\mathcal{G}_{\mathrm{A}}(a)\chi_{1} = \mathcal{G}_{\mathrm{A}}(a)\chi_{2}\)）。于是该图表蕴涵映射 \(\mathsf{X}(h)\) 是满射的。

\subsection{联合谱}

设 \(\Lambda\) 是一个集合。设 \(C_{\Lambda} = \mathbf{C}[(X_{\lambda})_{\lambda \in \Lambda}]\) 是不定元族 \((X_{\lambda})_{\lambda \in \Lambda}\) 上的复多项式所成的幺代数。对每个 \(\chi \in X(C_{\Lambda})\)，有 \((\chi (X_{\lambda}))_{\lambda \in \Lambda} \in C^{\Lambda}\)；映射 \(\chi \mapsto (\chi (X_{\lambda}))_{\lambda \in \Lambda}\) 是从 \(X(C_{\Lambda})\) 到乘积空间 \(C^{\Lambda}\) 上的一个同胚，借此把这两个空间视为同一。

另一方面，设 A 是一个交换幺 Banach 代数，并设 \(x = (x_{\lambda})_{\lambda \in \Lambda}\) 是 A 的元素的一个族。存在唯一的从 \(C_{\Lambda}\) 到 A 的保单位元态射 h，使得 \(h(\mathrm{X}_{\lambda}) = x_{\lambda}\) 对一切 \(\lambda\) 成立。连续映射 \(\mathsf{X}(h)\)（从 \(\mathsf{X}(\mathrm{A})\) 到 \(C^{\Lambda}\) 中）是把 \(\chi\) 联系到族 \((\chi(x_{\lambda}))_{\lambda \in \Lambda}\) 的映射。它称为由 x 定义的从 \(\mathsf{X}(\mathrm{A})\) 到 \(C^{\Lambda}\) 中的映射。

定义 2. --- 映射 X(h) 的像称为 x 的联合谱，记作 \(\mathrm{Sp}_{\mathrm{A}}^{\Lambda}(x)\) 或 \(\mathrm{Sp}^{\Lambda}(x)\)。

\(x\) 的联合谱是 \(\mathbf{C}^{\Lambda}\) 的一个紧子集。一个元素 \(c = (c_{\lambda}) \in \mathbf{C}^{\Lambda}\) 属于 \(\mathrm{Sp}_{\mathrm{A}}^{\Lambda}(x)\)，当且仅当诸元素 \(x_{\lambda} - c_{\lambda}\) 属于 \(\mathrm{A}\) 的同一个极大理想，换言之，当且仅当族 \((x_{\lambda} - c_{\lambda})_{\lambda \in \Lambda}\) 不生成代数 \(\mathrm{A}\)。

若 \(\Lambda\) 只含一个元素，从而族 \(x\) 归结为单个元素 \(x \in \mathsf{A}\)，则有 \(\mathrm{Sp}_{\mathsf{A}}^{\Lambda}(x) = \mathrm{Sp}_{\mathsf{A}}(x)\) (I, p.~37, 命题 6, \(b\))。若 \(\Lambda' \subset \Lambda\)，则 \(\mathrm{Sp}_{\mathsf{A}}^{\Lambda'}((x_{\lambda})_{\lambda \in \Lambda'})\) 是 \(\mathrm{Sp}_{\mathsf{A}}^{\Lambda}((x_{\lambda})_{\lambda \in \Lambda})\) 在从 \(\mathbf{C}^{\Lambda}\) 到 \(\mathbf{C}^{\Lambda'}\) 上的典范投影映射下的像。特别地，我们有

\[
\operatorname{Sp} _ {\mathrm{A}} ^ {\Lambda} (x) \subset \prod_ {\lambda \in \Lambda} \operatorname{Sp} _ {\mathrm{A}} (x _ {\lambda}).
\]

我们用 \(z_{\lambda}\)（\(\lambda \in \Lambda\)）表示 \(\mathbf{C}^{\Lambda}\) 上的坐标函数。若 \(\chi \in \mathsf{X}(\mathsf{A})\)，则在 \(\chi\) 处 \(z_{\lambda} \circ \mathsf{X}(h)\) 的值是 \(\chi(x_{\lambda})\)，因而 \(z_{\lambda} \circ \mathsf{X}(h) = \mathcal{G}(x_{\lambda})\)。

设 A 与 B 是交换幺 Banach 代数，\(\varphi\) 是从 A 到 B 的一个保单位元态射，而 \(x = (x_{\lambda})_{\lambda \in \Lambda}\) 是 A 的元素的一个族。设 \(\varphi(x)\) 为 B 的元素所成的族 \((\varphi(x_{\lambda}))_{\lambda \in \Lambda}\)。我们对每个 \(\chi \in \mathsf{X}(\mathsf{B})\) 及每个 \(\lambda \in \Lambda\) 有

\[
\chi (\varphi (x _ {\lambda})) = (\mathsf {X} (\varphi) (\chi)) (x _ {\lambda}),
\]

因而 \(\operatorname{Sp}_{\mathrm{B}}^{\Lambda}(\varphi(x)) \subset \operatorname{Sp}_{\mathrm{A}}^{\Lambda}(x)\)。图表

\[
\begin{array}{c} \mathsf {X} (\mathrm{B}) \xrightarrow {\mathsf {X} (\varphi)} \mathsf {X} (\mathrm{A}) \\ \Big \downarrow \qquad \qquad \qquad \qquad \Big \downarrow \\ \operatorname{Sp} _ {\mathrm{B}} ^ {\Lambda} (\varphi (x)) \xrightarrow {i} \operatorname{Sp} _ {\mathrm{A}} ^ {\Lambda} (x) \end{array}\tag{1}
\]

（其中 i 表示包含映射，竖直箭头表示由族 \(\varphi(x)\) 与 x 定义的映射）因而是交换的。

例. --- 设 \(K \subset C^{\Lambda}\) 是一个紧子集。设 \(z = (z_{\lambda})_{\lambda \in \Lambda}\) 是 \(\mathcal{C}(K)\) 中由 \(\mathbf{C}^{\Lambda}\) 的坐标函数在 K 上的限制所成的族。则联合谱 \(\mathrm{Sp}_{\mathcal{C}(\mathrm{K})}^{\Lambda}(z)\) 等于 K。事实上，依 I, p.~32, 推论 2，每个特征标 \(\chi\)（\(\mathcal{C}(\mathrm{K})\) 的特征标）都具有形式 \(f \mapsto f(x)\)（对某个元素 \(x \in \mathrm{K}\)），于是有 \((\chi(z_{\lambda}))_{\lambda \in \Lambda} = x\)。

命题 12. — 设 \(\Lambda\) 是一个集合，A 是一个交换幺 Banach 代数。设 \(x = (x_{\lambda})_{\lambda \in \Lambda}\) 是 A 的元素的一个族。

\begin{enumerate}
\def\labelenumi{\alph{enumi})}
\item
  假定由族 x 生成的 A 的全子代数在 A 中稠密。由 x 定义的从 X(A) 到 \(\mathbf{C}^{\Lambda}\) 的映射是从 X(A) 到联合谱 \(\mathrm{Sp}_{\mathrm{A}}^{\Lambda}(x)\) 上的一个同胚；
\item
  假定由族 x 生成的 A 的幺子代数在 A 中稠密。则对每个 \(c \in C^{\Lambda}\)，下列条件等价：
\end{enumerate}

\begin{enumerate}
\def\labelenumi{(\roman{enumi})}
\item
  \(c \in \mathrm{Sp}_{\mathrm{A}}^{\Lambda}(x)\)；
\item
  \(|\mathrm{P}(c)| \leqslant \varrho(\mathrm{P}(x))\) 对每个多项式 \(\mathrm{P} \in \mathbf{C}[(\mathrm{X}_{\lambda})_{\lambda \in \Lambda}]\) 成立；
\item
  \(|\mathrm{P}(c)| \leqslant \| \mathrm{P}(x)\|\) 对每个多项式 \(\mathrm{P} \in \mathbf{C}[(\mathrm{X}_{\lambda})_{\lambda \in \Lambda}]\) 成立。
\end{enumerate}

\begin{enumerate}
\def\labelenumi{\alph{enumi})}
\item
  从 \(\mathsf{X}(\mathsf{A})\) 到 \(\mathrm{Sp}_{\mathrm{A}}^{\Lambda}(x)\) 中的、由族 \(x\) 定义的映射是连续且满射的。设 \(\chi\), \(\chi' \in \mathsf{X}(\mathsf{A})\) 是具有同一个像的两个特征标，即 \(\chi(x_{\lambda}) = \chi'(x_{\lambda})\) 对一切 \(\lambda \in \Lambda\) 成立。特征标 \(\chi\) 与 \(\chi'\) 在形如 \(\mathrm{P}(x)\mathrm{Q}(x)^{-1}\) 的元素上一致，其中 \(\mathrm{P} \in \mathbf{C}[(\mathrm{X}_{\lambda})]\)，\(\mathrm{Q} \in \mathbf{C}[(\mathrm{X}_{\lambda})]\) 且 \(\mathrm{Q}(x)\) 在 A 中可逆，也就是说，它们在由诸元素 \(x_{\lambda}\) 生成的 A 的全子代数上一致 (I, p.~6, 引理 2)。由于 \(\chi\) 与 \(\chi'\) 连续 (I, p.~29, 定理 1)，它们因此在 A 上相等。这就证明 \(\mathsf{X}(h)\) 是从 \(\mathsf{X}(\mathsf{A})\) 到 \(\mathrm{Sp}_{\mathrm{A}}^{\Lambda}(x)\) 上的一个连续双射，从而由于 \(\mathsf{X}(\mathsf{A})\) 是紧的，它是一个同胚。
\item
  我们来证明 (i) 蕴涵 (ii)：若 \(c = (c_{\lambda})_{\lambda \in \Lambda} \in \mathrm{Sp}_{\mathrm{A}}^{\Lambda}(x)\)，则存在 \(\chi \in \mathsf{X}(\mathsf{A})\) 使得 \(c_{\lambda} = \chi(x_{\lambda})\) 对一切 \(\lambda\) 成立。于是对每个 \(\mathrm{P} \in \mathbf{C}[(\mathrm{X}_{\lambda})]\)，我们有 \(|\mathrm{P}(c)| = |\mathrm{P}((\chi(x_{\lambda}))_{\lambda \in \Lambda})| = |\chi(\mathrm{P}(x))| \leqslant \varrho(\mathrm{P}(x))\)。
\end{enumerate}

断言 (ii) 蕴涵 (iii) 是依据不等式 \(\varrho(x) \leqslant \|x\|\)，它按谱半径的定义对一切 \(x \in A\) 成立。

最后我们来证明 (iii) 蕴涵 (i)。设 \(c = (c_{\lambda})_{\lambda \in \Lambda} \in \mathbf{C}^{\Lambda}\) 满足

\[
| \mathrm{P} (c) | \leqslant \| \mathrm{P} (x) \|\tag{2}
\]

对一切 \(P \in \mathbf{C}[(X_{\lambda})]\) 成立。设 A' 是由族 x 生成的 A 的幺子代数；它的元素形如 \(P(x)\)，其中 \(P \in \mathbf{C}[(X_{\lambda})]\)。估计式 (2) 蕴涵：条件 \(P(x) = 0\) 蕴涵 \(\mathrm{P}(c)=0\)。于是存在一个幺代数态射 \(\xi\)（从 \(A'\) 到 C 中），使得 \(\xi(x_{\lambda})=c_{\lambda}\) 对一切 \(\lambda\in\Lambda\) 成立。由 (2)，态射 \(\xi\) 是连续的。因而它由连续性延拓为一个特征标 \(\chi\)（\(\overline{\mathrm{A}'}=\mathrm{A}\) 的特征标），后者满足 \(c=(\chi(x_{\lambda}))_{\lambda\in\Lambda}\in\mathrm{Sp}_{\mathrm{A}}^{\Lambda}(x)\)。证毕。

\subsection{多项式凸集}

定义 3. --- 设 \(\Lambda\) 是一个集合，V 是 \(\mathbf{C}^{\Lambda}\) 的一个子集。我们说 V 是多项式凸的，如果 V 是点 \((c_{\lambda})_{\lambda \in \Lambda}\)（\(\mathbf{C}^{\Lambda}\) 中的点）所成的集合，这些点满足

\[
| \mathrm{P} ((c _ {\lambda})) | \leqslant \sup _ {c \in \mathrm{V}} | \mathrm{P} (c) |
\]

对一切 \(\mathrm{P}\in \mathbf{C}[(\mathrm{X}_{\lambda})]\) 成立。

引理 3. --- 设 \(\Lambda\) 是一个集合。\(\mathbf{C}^{\Lambda}\) 的一个子集 V 是多项式凸的，当且仅当存在一个族 \((\mathrm{P}_i)_{i\in \mathrm{I}}\)（由 \(\mathbf{C}[(\mathrm{X}_{\lambda})_{\lambda \in \Lambda}]\) 的元素组成）与一个族 \((\mathrm{M}_i)_{i\in \mathrm{I}}\)（由 \([0, +\infty]\) 的元素组成），使得 V 是 \(c\in \mathbf{C}^{\Lambda}\) 中满足

\[
| \mathrm{P} _ {i} (c) | \leqslant \mathrm{M} _ {i}
\]

（对一切 i ∈ I）的元素所成的集合。

若 \(C^{\Lambda}\) 的子集 V 是多项式凸的，则它对由 \(\mathbf{C}[(X_{\lambda})_{\lambda\in\Lambda}]\) 的元素 P 组成的族满足上述条件，其中取 \(M_{P}=\sup_{c\in V}|P(c)|\)。

反之，设 \((\mathrm{P}_{i})_{i\in\mathrm{I}}\) 是由 \(\mathbf{C}[(\mathrm{X}_{\lambda})_{\lambda\in\Lambda}]\) 的元素组成的一个族，并设 \((\mathrm{M}_{i})_{i\in\mathrm{I}}\) 是由 \([0,+\infty]\) 的元素组成的一个族。设 V 是 \(C^{\Lambda}\) 中使得 \(|\mathrm{P}_{i}(c)|\leqslant\mathrm{M}_{i}\)（对一切 \(i\in I\)）的 c 所成的集合。于是我们有 \(\sup_{c\in\mathrm{V}}|\mathrm{P}_{i}(c)|\leqslant\mathrm{M}_{i}\) 对 \(i\in I\) 成立。假定 \(x\in C^{\Lambda}\) 满足

\[
| \mathrm{P} (x) | \leqslant \sup _ {c \in \mathrm{V}} | \mathrm{P} (c) |
\]

对一切 \(\mathrm{P}\in\mathbf{C}[(\mathrm{X}_{\lambda})]\) 成立。特别地，对 \(i\in I\)，我们有 \(|\mathrm{P}_{i}(x)|\leqslant\mathrm{M}_{i}\)，因而 \(x\in V\)。反过来，对每个元素 \(x\in V\) 及每个多项式 \(\mathrm{P}\in\mathbf{C}[(\mathrm{X}_{\lambda})]\)，有 \(|\mathrm{P}(x)|\leqslant\sup_{c\in\mathrm{V}}|\mathrm{P}(c)|\)。因此，集合 V 是多项式凸的。

引理 4. --- 设 A 是一个交换幺 Banach 代数。设 \(\Lambda\) 是一个集合，\(x = (x_{\lambda})_{\lambda \in \Lambda}\) 是 A 的元素的一个族。若由族 x 生成的幺子代数在 A 中稠密，则联合谱 \(\mathrm{Sp}_{\mathrm{A}}^{\Lambda}(x)\) 是多项式凸的。

这由 I, p.~43, 命题 12 的断言 \(b\) ) 以及定义 3 得出。

\(C^{\Lambda}\) 的多项式凸子集的任意交是多项式凸的（引理 3）。这证明了下述定义的合理性：

定义 4. --- 设 \(\Lambda\) 是一个集合，V 是 \(\mathbf{C}^{\Lambda}\) 的一个子集。V 的多项式凸包是 \(\mathbf{C}^{\Lambda}\) 中包含 V 的最小多项式凸子集。

V 的多项式凸包是 \(C^{\Lambda}\) 中满足 \(|P(c)| \leqslant \sup_{W} |P|\)（对一切 \(P \in \mathbf{C}[(X_{\lambda})]\)）的 c 所成的集合。事实上，依引理 3 这个集合是多项式凸的，且依定义它含于每个包含 V 的多项式凸集中。

例. --- 设 \(\Lambda\) 是一个有限集合，\(V \subset \mathbf{C}^{\Lambda}\) 是一个紧凸子集。则 \(V\) 是多项式凸的。事实上，设 \(W\) 是 \(V\) 的多项式凸包。我们来证明 \(W \subset V\)，由此即得该断言。设 \(x \in \mathbf{C}^{\Lambda} - V\)。存在一个实超平面 \(H\)（\(\mathbf{C}^{\Lambda}\) 中的实超平面），它严格分离 \(x\) 与 \(V\) (EVT, II, p.~41, 命题 4)。设 \(f_{\mathbf{R}}\) 是一个 \(\mathbf{R}\) -线性形式，定义在 \(\mathbf{C}^{\Lambda}\) 上，并设 \(\alpha \in \mathbf{R}\)，使得 \(H\) 是 \(y \in \mathbf{C}^{\Lambda}\) 中满足 \(f_{\mathbf{R}}(y) = \alpha\) 的 y 所成的集合。设 \(f\) 是 \(\mathbf{C}^{\Lambda}\) 上的一个线性形式，使得 \(f_{\mathbf{R}} = \mathcal{R}(f)\)。于是我们有

\[
\mathcal {R} (f (x)) > \sup _ {y \in V} \mathcal {R} (f (y)).
\]

对每个 \(t \in R\) 与 \(y \in V\)，令 \(f_{t}(y) = t + f(y)\)。我们有 \(|f_{t}| - \mathcal{R}(f_{t}) \to 0\)（在赋以紧收敛拓扑的 \(\mathcal{C}(\mathbf{C}^{\Lambda}, \mathbf{R})\) 中，当 \(t \to +\infty\) 时）。对充分大的 t，由于 V 是紧的，由此得出 \(|f_{t}(x)| > \sup_{y \in V} |f_{t}(y)|\)。于是 \(x \in C^{\Lambda} - W\)，因为 \(f_{t}\) 是多项式函数。

引理 5. --- 设 K 是 C 的一个紧子集。我们用 \(\widehat{\mathbf{K}}\) 表示 K 与 \(\mathbf{C} - \mathbf{K}\) 的相对紧连通分支的并。则集合 \(\widehat{\mathbf{K}}\) 是紧的。

由于 K 是紧的，存在实数 r > 0，使得 K 含于以 0 为中心、以 r 为半径的开圆盘 D。于是 C-K 包含 C-D。由于空间 C-D 是连通的（它同胚于 \([r, +\infty[\times\mathbf{S}^1)\)），它含于 \(\mathbf{C}-\mathbf{K}\) 的一个连通分支 U。\(\mathbf{C}-\mathbf{K}\) 的任何其他连通分支都含于 D，因而都是有界的。于是连通分支 U 是 \(\mathbf{C}-\mathbf{K}\) 的唯一的无界连通分支，即 \(\mathrm{U} = \mathbf{C}-\widehat{\mathrm{K}}\)。由于 U 是开的且包含圆盘 D 的余集，集合 \(\widehat{\mathrm{K}}\) 是紧的。

命题 13. --- 设 \(n \geqslant 1\) 是一个整数。设 \(\mathbf{K} \subset \mathbf{C}^{n}\) 是一个闭子集，V 是它的多项式凸包。

\begin{enumerate}
\def\labelenumi{\alph{enumi})}
\item
  \(C^{n} - K\) 的每个有界连通分支都含于 V；
\item
  若 n = 1 且 K 是紧的，则 V 是 K 与 C-K 的有界连通分支的并 \(\widehat{K}\)。
\end{enumerate}

由于 \(K \subset C^{n}\) 是闭的，它的余集 \(C^{n}-K\) 是开的，从而是局部连通的，因此 \(C^{n}-K\) 的每个连通分支都是开的。最大模原理 (VAR, R1, p.~29, 3.3.7) 表明 \(C^{n}-K\) 的每个有界连通分支都含于 V，这就证明了断言 a)。

现在假定 \(n = 1\) 且 K 是紧的。用 \(\widehat{\mathbf{K}}\) 表示 K 与 \(\mathbf{C} - \mathbf{K}\) 的有界连通分支的并，于是 \(\widehat{\mathbf{K}} \subset \mathbf{V}\)；由前所述，集合 \(\widehat{\mathbf{K}}\) 是紧的（引理 5）。设 A 是交换幺 Banach 代数 \(\mathcal{C}(\widehat{\mathbf{K}})\)。设 \(x \in \mathbf{A}\) 是 \(\widehat{\mathbf{K}}\) 上的恒等函数，B 是 A 的由 \(x\) 生成的闭幺子代数。我们有 \(\mathrm{Sp}_{\mathrm{A}}(x) = \widehat{\mathbf{K}}\) (I, p.~17, 例 3)，因而 \(\mathbf{C} - \mathrm{Sp}_{\mathrm{A}}(x)\) 是连通的。于是我们有 \(\mathrm{Sp}_{\mathrm{B}}(x) = \mathrm{Sp}_{\mathrm{A}}(x)\) (I, p.~28, 命题 6 的推论)。由于依引理 4 \(\mathrm{Sp}_{\mathrm{B}}(x)\) 是多项式凸的，这表明 \(\widehat{\mathbf{K}}\) 是多项式凸的，从而 \(\mathbf{V} \subset \widehat{\mathbf{K}}\)。

该命题的第二部分不能推广到 \(n \geqslant 2\) 的情形（参见 I, p.~170, 习题 23）。

命题 14. --- 设 \(\Lambda\) 是一个集合。设 A 是一个交换幺 Banach 代数，\(x = (x_{\lambda})_{\lambda \in \Lambda}\) 是 A 的元素的一个族，A' 是由 \(x\) 生成的幺 Banach 子代数。则联合谱 \(\mathrm{Sp}_{\mathrm{A}'}^{\Lambda}(x)\) 是 \(\mathrm{Sp}_{\mathrm{A}}^{\Lambda}(x)\) 的多项式凸包。

事实上，I, p.~43, 命题 12 b) 证明了 \(\mathrm{Sp}_{\mathrm{A}'}^{\Lambda}(x)\) 是由 \(c \in \mathbf{C}^{\Lambda}\) 中满足 \(|\mathrm{P}(c)| \leqslant \varrho(\mathrm{P}(x))\)（对一切 \(\mathrm{P} \in \mathbf{C}[(\mathrm{X}_{\lambda})]\)）的 c 所成的集合。

现在我们有

\[
\varrho (\mathrm{P} (x)) = \sup _ {\chi \in \mathsf {X} (\mathrm{A})} | \chi (\mathrm{P} (x)) | = \sup _ {\chi \in \mathsf {X} (\mathrm{A})} | \mathrm{P} ((\chi (x _ {\lambda})) _ {\lambda \in \Lambda}) | = \sup _ {c \in \operatorname{Sp} _ {\mathrm{A}} ^ {\Lambda} (x)} | \mathrm{P} (c) |,
\]

这是依 I, p.~38, 命题 7 a) 及 I, p.~43, 命题 12 a)。于是由引理 3 即得结论。

命题 15. --- 设 \(\Lambda\) 是一个集合。设 K 是 \(\mathbf{C}^{\Lambda}\) 的一个紧多项式凸子集。设 \(\mathrm{A}_0\) 是 \(\mathbf{C}^{\Lambda}\) 上的多项式函数在 K 上的限制所成的集合，并设 A 是 \(\mathrm{A}_0\) 在代数 \(\mathcal{C}(\mathrm{K})\) 中的闭包。

设 ev: K → X(A) 是由 x ↦ ev\_x 定义的映射，其中 ev\_x 是 A 的特征标 f ↦ f(x)。设 z = (z\_λ)\_\{λ∈Λ\} 是 A 中由 C\^{}Λ 上的坐标函数在 K 上的限制所成的族，并设 φ 是由族 z 定义的、从 X(A) 到 Sp\_A\^{}Λ(z) 上的满射。

于是我们有 K = Sp \(_{A}^{\Lambda}(z)\)，且映射 ev: K → X(A) 与 φ: X(A) → K 是互逆的同胚。

映射 \(\varphi \circ\) ev 是 K 的恒等映射。特别地，K 含于像 \(\mathrm{Sp}_{\mathrm{A}}^{\Lambda}(z)\)（\(\varphi\) 的像）之中。

命题 14 蕴涵 \(\mathrm{Sp}_{\mathrm{A}}^{\Lambda}(z)\) 是 \(\mathrm{Sp}_{\mathcal{C}(\mathrm{K})}^{\Lambda}(z)\) 的多项式凸包。另一方面，由于我们有 \(\mathrm{Sp}_{\mathcal{C}(\mathrm{K})}^{\Lambda}(z) = \mathrm{K}\) (I, p.~42, 例)，且依假设它是多项式凸的，我们推出 \(\mathrm{Sp}_{\mathrm{A}}^{\Lambda}(z) = \mathrm{K}\)。

由于由元素族 \(z_{\lambda}\) 生成的幺子代数在 A 中稠密，映射 \(\varphi\) 是从 \(\mathsf{X}(\mathsf{A})\) 到 \(\mathrm{Sp}_{\mathsf{A}}^{\Lambda}(z)=\mathrm{K}\) 上的一个同胚（第 I 卷, p. 43, 命题 12(a)）。恒等式 \(\varphi\circ ev=Id_{K}\) 于是表明 ev 是 \(\varphi\) 的逆同胚。

对每个子集 \(\Lambda'\)（\(\Lambda\) 的子集），我们用 \(\mathrm{pr}_{\Lambda'}\) 表示典范投影 \(\mathbf{C}^{\Lambda} \to \mathbf{C}^{\Lambda'}\)。设 W 是 \(\mathbf{C}^{\Lambda}\) 的一个子集，V 是它的多项式凸包。令 \(W' = \mathrm{pr}_{\Lambda'} W\)。由于 \(\mathbf{C}[(X_{\lambda})_{\lambda \in \Lambda'}]\) 的每个元素都与 \(\mathbf{C}[(X_{\lambda})_{\lambda \in \Lambda}]\) 的一个元素视为同一，\(W'\) 的多项式凸包含于 \(\mathrm{pr}_{\Lambda'} V\)。

引理 6. --- 设 \(K \subset C^{\Lambda}\) 是一个紧多项式凸子集，U 是 K 的一个邻域。存在一个有限子集 \(\Lambda_{0}\)（\(\Lambda\) 的子集），使得对每个子集 \(\Lambda'\)（\(\Lambda\) 的包含 \(\Lambda_{0}\) 的子集），集合 \(\mathrm{pr}_{\Lambda'}(\mathrm{U})\) 包含 \(\mathrm{pr}_{\Lambda'}(\mathrm{K})\) 的多项式凸包。

由于 K 是紧的，存在 C 中一族紧圆盘 \(D_{\lambda}\)（以 0 为中心，半径为 \(R_{\lambda}\)），使得 K 含于乘积 \(\mathrm{D} = \prod_{\lambda} \mathrm{D}_{\lambda}\)。对每个 \(\mathrm{P} \in \mathbf{C}[(\mathrm{X}_{\lambda})]\)，用 \(\mathrm{K}_{\mathrm{P}}\) 表示 \(x \in \mathbf{C}^{\Lambda}\) 中满足

\[
| \mathrm{P} (x) | \leqslant \sup _ {c \in \mathrm{K}} | \mathrm{P} (c) |.
\]

我们有

\[
\mathrm{D} \cap \bigcap_ {\mathrm{P}} \mathrm{K} _ {\mathrm{P}} = \mathrm{K}\tag{3}
\]

它含于 U。于是，空间 D 是开集 \(D \cap U\) 与诸开集 \(D - K_{P}\)（\(P \in \mathbf{C}[(X_{\lambda})]\)）的并。由于 D 是紧的，存在一个整数 \(q \geqslant 1\) 及多项式 \(P_{1}, \ldots, P_{q} \in \mathbf{C}[(X_{\lambda})]\)，使得：

\[
\mathrm{D} \cap \mathrm{K} _ {\mathrm{P} _ {1}} \cap \dots \cap \mathrm{K} _ {\mathrm{P} _ {q}} \subset \mathrm{U}.\tag{4}
\]

存在一个有限集 \(\Lambda_0 \subset \Lambda\)，使得 \(P_i \in \mathbf{C}[(X_\lambda)_{\lambda \in \Lambda_0}]\) 对 \(1 \leqslant i \leqslant q\) 成立。我们来证明 \(\Lambda_0\) 具有所要求的性质。

设 \(\Lambda'\) 是 \(\Lambda\) 的一个包含 \(\Lambda_0\) 的子集。设 E 是 \(\mathbf{C}^{\Lambda'}\) 的子集，由元素 \(c = (c_{\lambda})_{\lambda \in \Lambda'}\)（由不等式 \(|c_{\lambda}| \leqslant \mathrm{R}_{\lambda}\)（对 \(\lambda \in \Lambda'\)）所定义）组成，且

\[
| \mathrm{P} _ {i} (c) | \leqslant \sup _ {c \in \mathrm{K}} | \mathrm{P} _ {i} (c) |
\]

（对 \(i = 1, \ldots, q\)）。集合 E 是多项式凸的（引理 3），且公式 (3) 表明 \(\mathrm{pr}_{\Lambda'}(\mathrm{K}) \subset \mathrm{E}\)。

另一方面，设 \(c = (c_{\lambda})_{\lambda \in \Lambda'} \in \mathrm{E}\)；设 \(d = (d_{\lambda})_{\lambda \in \Lambda}\) 是 \(\mathbf{C}^{\Lambda}\) 中由 \(d_{\lambda} = c_{\lambda}\)（对 \(\lambda \in \Lambda'\)）与 \(d_{\lambda} = 0\)（对 \(\lambda \in \Lambda - \Lambda'\)）定义的元素。于是 (4) 蕴涵 \(d \in \mathrm{U}\)，从而 \(c \in \mathrm{pr}_{\Lambda'}(\mathrm{U})\)。这样，\(\mathrm{E} \subset \mathrm{pr}_{\Lambda'}(\mathrm{U})\)，证明完成。

引理 7. --- 设 \(n \geqslant 1\) 是一个整数，K 是 \(\mathbf{C}^{n}\) 的一个多项式凸紧子集。则 K 具有一个由多项式凸紧邻域组成的基本系。

存在一个紧多圆柱 (cf.~VAR, R1, p.~24) \(\Delta\)（\(\mathbf{C}^n\) 中的），它是 K 的一个邻域。由于 K 是多项式凸的，存在一个族 \((\mathrm{P}_i)_{i\in \mathrm{I}}\)（由 \(\mathbf{C}[\mathrm{X}_1,\dots ,\mathrm{X}_n]\) 的元素组成），以及一个由正实数组成的族 \((\mathrm{M}_i)_{i\in \mathrm{I}}\)，使得 K 是 \(z\in\) \(\Delta\) 中满足 \(|\mathrm{P}_i(z)|\leqslant \mathrm{M}_i\)（对一切 \(i\)）的 z 所成的集合（引理 3）。对 I 的每个有限子集 J 及每个 \(\varepsilon >0\)，用 \(\mathrm{K}_{\mathrm{J},\varepsilon}\) 表示 \(z\in \Delta\) 中满足 \(|\mathrm{P}_i(z)|\leqslant \mathrm{M}_i + \varepsilon\)（对 \(i\in \mathrm{J}\)）的 z 所成的集合。则每个集合 \(\mathrm{K}_{\mathrm{J},\varepsilon}\) 都是 K 的一个多项式凸紧邻域 (loc. cit.)，且诸集合 \(\mathrm{K}_{\mathrm{J},\varepsilon}\) 的交是 K。于是诸集合 \(\mathrm{K}_{\mathrm{J},\varepsilon}\) 构成 K 的多项式凸邻域的一个基本系 (TG, I, p.~60, th. 1)。

\section{全纯函数演算}

\subsection{全纯函数的芽}

设 E 与 F 是复 Banach 空间。我们回顾（cf. VAR, R1, p.~26, 3.2.1, p.~22, 3.1 及 p.~88, App.）：定义在 E 的开子集 U 上、取值于 F 中的全纯映射是这样的映射 \(f \colon \mathrm{U} \to \mathrm{F}\)：对一切 \(x \in \mathrm{U}\)，存在一个收敛级数

\[
f _ {x} = \sum_ {k \geqslant 0} f _ {x, k}
\]

它满足 \(f(x + y) = f_x(y)\)（对充分接近 0 的一切 \(y \in E\)），其中 \(f_{x,k} \colon E \to C\) 是 E 上的、取值于 F 中的 k 次连续齐次多项式，即形如

\[
f _ {x, k} (y) = \widetilde {f} _ {x, k} (y, \ldots , y)
\]

其中 \(\widetilde{f}_{x,k}\colon\mathrm{E}^{k}\to\mathrm{F}\) 是一个 \(k\) 线性连续映射。我们用 \(\mathcal{O}(\mathrm{U};\mathrm{F})\) 表示 U 上取值于 F 中的全纯函数所成的复向量空间，赋以紧收敛拓扑。它是一个局部凸拓扑向量空间，其拓扑由半范数 \(f\mapsto\sup_{z\in\mathrm{K}}\|f(z)\|\) 定义，这里 K 取遍 U 的紧子集所成的集合。

设 G 是一个复 Banach 空间，V 是 G 的一个开子集。对每个全纯映射 \(\varphi\colon\mathrm{V}\to\mathrm{U}\)，映射 \(\varphi^{*}\colon f\mapsto f\circ\varphi\) 是从 \(\mathcal{O}(\mathrm{U};\mathrm{F})\) 到 \(\mathcal{O}(\mathrm{V};\mathrm{F})\) 中的连续线性映射。

若 H 是一个复 Banach 空间，\(\varphi\colon\mathrm{F}\to\mathrm{H}\) 是一个连续线性映射，则映射 \(f\mapsto\varphi\circ f\) 是从 \(\mathcal{O}(\mathrm{U};\mathrm{F})\) 到 \(\mathcal{O}(\mathrm{U};\mathrm{H})\) 中的连续线性映射，记作 \(\varphi_{*}\)。

设 n 是一个自然数，并令 \(E = C^{n}\)。设 K 是 \(C^{n}\) 的一个紧子集，U 是 K 的开邻域组成的递降过滤集。若 U、\(U' \in U\) 且 \(U' \subset U\)，则从 \(\mathcal{O}(U; F)\) 到 \(\mathcal{O}(U'; F)\) 中的函数限制映射是连续的。诸空间 \(\mathcal{O}(\mathrm{U};\mathrm{F})\) 关于这些映射的归纳极限记作 \(\mathcal{O}(\mathrm{K};\mathrm{F})\)。\(\mathcal{O}(\mathrm{K};\mathrm{F})\) 的元素称为 K 的邻域中的、取值于 F 的全纯函数芽。

空间 \(\mathcal{O}(\mathrm{K};\mathrm{F})\) 赋以诸 \(\mathcal{O}(\mathrm{U};\mathrm{F})\) 的局部凸拓扑的归纳极限拓扑。设 X 是一个局部凸拓扑向量空间，\(\varphi\colon\mathcal{O}(\mathrm{K};\mathrm{F})\to\mathrm{X}\) 是一个映射。对 K 的每个开邻域 U，映射 \(\mathcal{O}(\mathrm{U};\mathrm{F})\to\mathrm{X}\)（由 \(\varphi\) 与典范映射 \(\mathcal{O}(\mathrm{U};\mathrm{F})\to\mathcal{O}(\mathrm{K};\mathrm{F})\) 复合而得到）记作 \(\varphi^{\mathrm{U}}\)。映射 \(\varphi\) 连续当且仅当对每个 U，\(\varphi^{\mathrm{U}}\) 连续 (EVT, II, p.~29, prop. 5, (iii))。

设 m 是一个自然数。设 L 是 \(C^{m}\) 的一个紧子集，V 是 L 的一个开邻域。设 \(\varphi\colon V \to C^{n}\) 是一个满足 \(\varphi(L) \subset K\) 的全纯映射。连续线性映射

\[
\mathcal {O} (\mathrm{U}; \mathrm{F}) \xrightarrow {\varphi^ {*}} \mathcal {O} (\stackrel {- 1} {\varphi} (\mathrm{U}); \mathrm{F}) \xrightarrow {\varphi^ {- 1} \stackrel {- 1} {\varphi} (\mathrm{U})} \mathcal {O} (\mathrm{L}; \mathrm{F})
\]

对 K 的每个开邻域 U，诱导出一个连续线性映射 \(\varphi^{*}\colon\mathcal{O}(\mathrm{K};\mathrm{F})\to\mathcal{O}(\mathrm{L};\mathrm{F})\) (loc. cit.)。

设 H 是一个复 Banach 空间，\(\varphi\colon\mathrm{F}\to\mathrm{H}\) 是一个连续线性映射。连续线性映射

\[
\mathcal {O} (\mathrm{U}; \mathrm{F}) \xrightarrow {\varphi_ {*}} \mathcal {O} (\mathrm{U}; \mathrm{H}) \xrightarrow {\varphi^ {\mathrm{U}}} \mathcal {O} (\mathrm{K}; \mathrm{F})
\]

（其中 U 取遍 K 在 \(C^{n}\) 中的开邻域所成的集合）诱导出一个连续线性映射 \(\varphi_{*}\)（从 \(\mathcal{O}(K;F)\) 到 \(\mathcal{O}(K;H)\) 中）(loc. cit.)。我们有时将记 \(\varphi \circ f = \varphi_{*}(f)\)。

对 K 的每个开邻域 U，在 K 上的限制是一个连续线性映射 \(\mathcal{O}(\mathrm{U};\mathrm{F})\to \mathcal{C}(\mathrm{K};\mathrm{F})\)；这些映射诱导出一个连续线性映射 \(\mathcal{O}(\mathrm{K};\mathrm{F})\to \mathcal{C}(\mathrm{K};\mathrm{F})\)，称为 K 上全纯函数芽的赋值。

设 A 是一个复幺 Banach 代数。空间 \(\mathcal{O}(\mathrm{U};\mathrm{A})\) 与 \(\mathcal{O}(\mathrm{K};\mathrm{A})\) 是幺代数。若 \(A \neq \{0\}\)，则可以把 \(\mathcal{O}(\mathrm{U};\mathbf{C})\)（相应地，\(\mathcal{O}(\mathrm{K};\mathbf{C})\)）典范地视为子代数 \(\mathcal{O}(\mathrm{U};\mathbf{C}) \cdot 1\)（\(\mathcal{O}(\mathrm{U};\mathrm{A})\) 的子代数）（相应地，子代数 \(\mathcal{O}(\mathrm{K};\mathbf{C}) \cdot 1\)（\(\mathcal{O}(\mathrm{K};\mathrm{A})\) 的子代数））。我们将令 \(\mathcal{O}(\mathrm{U}) = \mathcal{O}(\mathrm{U};\mathbf{C})\) 及 \(\mathcal{O}(\mathrm{K}) = \mathcal{O}(\mathrm{K};\mathbf{C})\)。

\subsection{主要定理的叙述}

设 X 是一个集合。若 \(m \leqslant n\)，我们将用 \(\pi_{m,n}\) 表示从 \(X^{n}\) 到 \(X^{m}\) 中的、使得 \(\pi_{m,n}(\boldsymbol{x}) = (x_{1}, \ldots, x_{m})\)（对一切 \(\boldsymbol{x} = (x_{1}, \ldots, x_{n}) \in \mathrm{X}^{n}\)）的映射。

设 A 是 C 上的一个幺 Banach 代数。对每个整数 \(n \geqslant 1\) 及每个 \(a \in A^{n}\)，我们用 \(\mathrm{Sp}^{n}(\boldsymbol{a})\) 表示联合谱 \(\mathrm{Sp}_{\mathrm{A}}^{\{1,\ldots,n\}}(\boldsymbol{a})\) (I, p.~42, 定义 2)。它是 \(C^{n}\) 的一个紧子集。对每个满足 \(1 \leqslant m \leqslant n\) 的整数 m，有 \(\pi_{m,n}(\mathrm{Sp}^{n}(\boldsymbol{a})) = \mathrm{Sp}^{m}(\pi_{m,n}(\boldsymbol{a}))\) (I, p.~41, 第 6 小节)。连续线性映射

\[
\pi_ {m, n} ^ {*} \colon \mathcal {O} (\mathrm{Sp} ^ {m} (\pi_ {m, n} (\boldsymbol {a})); \mathrm{A}) \longrightarrow \mathcal {O} (\mathrm{Sp} ^ {n} (\boldsymbol {a}); \mathrm{A})
\]

是幺代数的一个态射。

设 A 是 C 上的一个交换幺 Banach 代数。设 \(n \geqslant 1\) 是一个整数。我们把这样的资料称为 A 上 n 个变量的全纯函数演算：对每个 \(a \in A^{n}\)，给定一个映射

\[
\Theta_ {\boldsymbol {a}} \colon \mathcal {O} (\mathrm{Sp} ^ {n} (\boldsymbol {a}); \mathrm{A}) \longrightarrow \mathrm{A}
\]

满足下列条件：

(CF1) 对每个 \(a \in A^{n}\)，映射 \(\Theta_{a}\) 是幺代数的一个连续态射。

(CF2) 若 \(\boldsymbol{a}=(a_{1},\ldots,a_{n})\)，且 \(z_{1},\ldots,z_{n}\) 表示 \(\mathrm{Sp}^{n}(\boldsymbol{a})\) 的邻域中 \(C^{n}\) 的坐标函数的芽，则有

\[
\Theta_ {\boldsymbol {a}} (z _ {1}) = a _ {1}, \dots , \Theta_ {\boldsymbol {a}} (z _ {n}) = a _ {n}.
\]

注. --- 若代数 A 的根为零，则可以省去 (CF1) 中的连续性条件（参见 I, p.~40, 命题 9）。

我们把这样的资料称为 A 上的全纯函数演算：对每个整数 \(n \geqslant 1\)，给定 A 上 n 个变量的一个全纯函数演算，且满足：

(CF3) 无论 m 与 n 是怎样的整数，只要 \(1 \leqslant m \leqslant n\)，也无论 \(a \in A^{n}\) 与 \(f \in \mathcal{O}(\mathrm{Sp}^{m}(\pi_{m,n}(a); \mathrm{A})\) 是怎样的，总有

\[
\Theta_ {\boldsymbol {a}} (\pi_ {m, n} ^ {*} (f)) = \Theta_ {\pi_ {m, n} (\boldsymbol {a})} (f).
\]

本节的目的在于证明下述定理：

定理 1. --- 设 A 是一个复交换幺 Banach 代数。则存在唯一的一个 A 上的全纯函数演算。

这一定理的证明将占据第 3 至第 7 小节。

\subsection{适配序列与相伴微分形式}

在本小节中，直到第 5 小节，我们用 A 表示一个复交换幺 Banach 代数，用 n 表示一个整数 \(\geqslant\) 1。

当我们谈到 \(C^{n}\) 的开子集上的无穷次可微函数时，指的是对于底实流形结构无穷次可微的函数。所用的微分演算概念都是相对于这个结构的。

定义 1. --- 设 \(\boldsymbol{a} = (a_{1},\ldots ,a_{n})\in \mathrm{A}^{n}\)，设 \(h\colon \mathbf{C}^n\to \mathbf{C}\) 是一个映射，并设 \(u_{1},\ldots ,u_{n}\) 是从 \(\mathbf{C}^n\) 到 A 中的映射。我们说序列 \((h,u_1,\dots ,u_n)\) 适配于 \(\boldsymbol{a}\)，如果

\begin{enumerate}
\def\labelenumi{(\roman{enumi})}
\item
  映射 h 是无穷次可微的，具有紧支集，且在 \(\mathrm{Sp}^n (\boldsymbol {a})\) 的一个邻域中等于 1；
\item
  诸映射 \(u_1, \ldots, u_n\) 都是无穷次可微的；
\item
  对每个 \(\boldsymbol{z} = (z_{1},\dots ,z_{n})\in \mathbf{C}^{n}\)，我们有
\end{enumerate}

\[
h (\boldsymbol {z}) + (z _ {1} - a _ {1}) u _ {1} (\boldsymbol {z}) + \dots + (z _ {n} - a _ {n}) u _ {n} (\boldsymbol {z}) = 1.\tag{1}
\]

\(C^{n}\) 上的、以 A 为系数的 2n 次微分形式，由

\[
\omega = \bigwedge_ {i = 1} ^ {n} (d u _ {i} \wedge d z _ {i})
\]

所定义，称为与 \((h, u_1, \ldots, u_n)\) 相伴的微分形式。

若 \((h, u_{1}, \ldots, u_{n})\) 适配于 a，则对 (1) 微分，我们得到等式

\[
d h = - \sum_ {i = 1} ^ {n} u _ {i} d z _ {i} - \sum_ {i = 1} ^ {n} (z _ {i} - a _ {i}) d u _ {i}\tag{2}
\]

由此，对每个满足 \(1 \leqslant i \leqslant n\) 的 i，得到关系式

\[
d h \wedge d z _ {i} \wedge \bigwedge_ {j \neq i} (d u _ {j} \wedge d z _ {j}) = - (z _ {i} - a _ {i}) \omega .\tag{3}
\]

引理 1. --- 设 U 是 \(\mathbf{C}^{n}\) 的一个开子集。设 K 是 U 的一个紧子集。则存在一个从 \(\mathbf{C}^{n}\) 到 C 中的无穷次可微函数 h，它在 K 上等于 1 且具有含于 U 中的紧支集。

设 V 是 K 的一个相对紧开邻域，使得 \(\overline{V}\) 含于 U (TG, I, p.~65, prop. 10)。存在 \(C^{n}\) 到 C 中的一个无穷次可微函数 h，其支集含于 V 且在 K 上等于 1 (VAR, R1, p.~40, 5.3.6)。这个函数具有所要求的性质。

例. --- 假定 n = 1。设 \(a \in A\)。对 \(\mathrm{Sp}(a)\) 的每个开邻域 U，存在一个从 C 到 C 中的无穷次可微函数 h，其紧支集含于 U，且在 \(\mathrm{Sp}(a)\) 的一个邻域中等于 1 (VAR, R1, p.~40, 5.3.6)。令

\[
u (z) = (1 - h (z)) (z - a) ^ {- 1}
\]

（对 \(z \in \mathbf{C} - \mathrm{Sp}(a)\)），并令 \(u(z) = 0\)（若 \(z \in \mathrm{Sp}(a)\)）。数对 \((h, u)\) 适配于 a，而相伴微分形式是 \(\omega = du \wedge dz\)。

引理 2. --- 设 \(\boldsymbol{a} = (a_{1},\ldots ,a_{n})\in \mathrm{A}^{n}\)。存在无穷次可微映射 \(v_{1},\ldots ,v_{n}\)（从 \(\mathbf{C}^n -\mathrm{Sp}^n (\boldsymbol {a})\) 到 A 中），使得

\[
(z _ {1} - a _ {1}) v _ {1} (\boldsymbol {z}) + \dots + (z _ {n} - a _ {n}) v _ {n} (\boldsymbol {z}) = 1
\]

对一切 \(z = (z_{1}, \ldots, z_{n}) \in \mathbf{C}^{n} - \mathrm{Sp}^{n}(\boldsymbol{a})\)。

设 \(\boldsymbol{w} = (w_{1}, \ldots, w_{n}) \in \mathbf{C}^{n} - \operatorname{Sp}^{n}(\boldsymbol{a})\)。依联合谱的定义，存在 A 中的元素 \(b_{1}, \ldots, b_{n}\)，使得

\[
(w _ {1} - a _ {1}) b _ {1} + \dots + (w _ {n} - a _ {n}) b _ {n} = 1
\]

(I, p.~41, 第 6 小节)。存在 w 的一个开邻域 \(W_{w}\)，使得 A 的元素 \((z_{1}-a_{1})b_{1}+\cdots+(z_{n}-a_{n})b_{n}\) 当 \(z=(z_{1},\ldots,z_{n})\) 属于 \(W_{w}\) 时是可逆的。于是，对每个满足 \(1\leqslant j\leqslant n\) 的整数 j 及 \(W_{w}\) 中的每个 z，令

\[
u _ {j} (\boldsymbol {z}) = b _ {j} \left(\sum_ {i = 1} ^ {n} (z _ {i} - a _ {i}) b _ {i}\right) ^ {- 1}.
\]

这样定义的函数 \(u_{1}, u_{2}, \ldots, u_{n}\)（\(W_{w}\) 到 A 中的函数）在 \(W_{w}\) 中是无穷次可微的，且我们有

\[
(z _ {1} - a _ {1}) u _ {1} (\boldsymbol {z}) + \dots + (z _ {n} - a _ {n}) u _ {n} (\boldsymbol {z}) = 1
\]

对 \(W_{w}\) 中的一切 z。

由于族 \((\mathrm{W}_{\boldsymbol{w}})_{\boldsymbol{w}\in\mathbf{C}^{n}-\mathrm{Sp}^{n}(\boldsymbol{a})}\) 是 \(\mathbf{C}^{n}-\mathrm{Sp}^{n}(\boldsymbol{a})\) 的一个开覆盖，存在一个局部有限的开加细 \(\mathscr{W}=(\mathrm{W}_{\lambda})_{\lambda\in\mathrm{L}}\) (TG, I, p.~70, th. 5)，且对每个 \(\lambda\in L\)，存在函数 \(u_{1\lambda},\ldots,u_{n\lambda}\)（取值于 A 中，定义且无穷次可微于 \(W_{\lambda}\) 中），使得 \((z_{1}-a_{1})u_{1\lambda}(\boldsymbol{z})+\cdots+(z_{n}-a_{n})u_{n\lambda}(\boldsymbol{z})=1\) 对一切 \(\boldsymbol{z}\)（位于 \(W_{\lambda}\) 中）成立。设 \((f_{\lambda})_{\lambda \in L}\) 是一个从属于覆盖 \(\mathscr{W}\) 的、由无穷次可微函数组成的单位分解 (VAR, R1, p.~40, 5.3.6)。设 \(i\) 是一个满足 \(1 \leqslant i \leqslant n\) 的整数。对每个 \(\lambda \in L\)，设 \(u_{i\lambda}'\) 是从 \(\mathbf{C}^n - \mathrm{Sp}^n(\boldsymbol{a})\) 到 A 中的映射，它由把函数 \(f_{\lambda}u_{i\lambda}\) 延拓（在 \((\mathbf{C}^n - \mathrm{Sp}^n(\boldsymbol{a})) - W_{\lambda}\) 中取 0）而得到。诸函数 \(u_{i\lambda}'\) 是无穷次可微的。由于族 \((\mathrm{Supp}(u_{i\lambda}')_{\lambda \in L}\) 是局部有限的，函数 \(v_i = \sum_{\lambda \in L} u_{i\lambda}'\) 在 \(\mathbf{C}^n - \mathrm{Sp}^n(\boldsymbol{a})\) 中有定义且无穷次可微。

设 \(z \in \mathbf{C}^n - \mathrm{Sp}^n(\boldsymbol{a})\)。用 L' 表示 \(\lambda \in \mathrm{L}\)（满足 \(z \in \mathrm{W}_{\lambda}\)）所成的有限集。则

\[
\begin{array}{c} \sum_ {i = 1} ^ {n} (z _ {i} - a _ {i}) v _ {i} (\boldsymbol {z}) = \sum_ {\lambda \in \mathrm{L} ^ {\prime}} \sum_ {i = 1} ^ {n} (z _ {i} - a _ {i}) u _ {i \lambda} ^ {\prime} (\boldsymbol {z}) \\ = \sum_ {\lambda \in \mathrm{L} ^ {\prime}} f _ {\lambda} (\boldsymbol {z}) \sum_ {i = 1} ^ {n} (z _ {i} - a _ {i}) u _ {i \lambda} (\boldsymbol {z}) = \bigg (\sum_ {\lambda \in \mathrm{L} ^ {\prime}} f _ {\lambda} (\boldsymbol {z}) \bigg) \cdot 1 = 1. \end{array}
\]

引理 3. --- 设 \(\boldsymbol{a} = (a_{1},\ldots ,a_{n})\in \mathrm{A}^{n}\)。设 \(h\) 是一个从 \(\mathbf{C}^n\) 到 \(\mathbf{C}\) 中的无穷次可微映射，在 \(\mathrm{Sp}^n (\boldsymbol {a})\) 的一个邻域中等于 1 且具有紧支集。存在无穷次可微映射 \(u_{1},\ldots ,u_{n}\)（从 \(\mathbf{C}^n\) 到 A 中的），使得序列 \((h,u_1,\dots,u_n)\) 适配于 \(\boldsymbol{a}\)。

设 \(v_{1},\ldots,v_{n}\) 是从 \(\mathbf{C}^{n}-\mathrm{Sp}^{n}(\boldsymbol{a})\) 到 A 中的无穷次可微映射，使得

\[
\sum_ {j = 1} ^ {n} (z _ {j} - a _ {j}) v _ {j} (\boldsymbol {z}) = 1
\]

（对 \(z\) 属于 \(\mathbf{C}^n - \mathrm{Sp}^n(\boldsymbol{a})\)）（引理 2）。设 \(i\) 是一个满足 \(1 \leqslant i \leqslant n\) 的整数。令 \(u_i(z) = (1 - h(z))v_i(z)\)（若 \(z \in \mathbf{C}^n - \mathrm{Sp}^n(\boldsymbol{a})\)），并令 \(u_i(z) = 0\)（若 \(z \in \mathrm{Sp}^n(\boldsymbol{a})\)）。诸映射 \(u_i\) 在 \(\mathbf{C}^n - \mathrm{Sp}^n(\boldsymbol{a})\) 中无穷次可微，且在 \(\mathrm{Sp}^n(\boldsymbol{a})\) 的一个邻域中为零，因而在 \(\mathbf{C}^n\) 中无穷次可微。等式 (1) 在 \(\mathrm{Sp}^n(\boldsymbol{a})\) 上成立，因为诸函数 \(u_i\) 在 \(\mathrm{Sp}^n(\boldsymbol{a})\) 上为零，且 \(h\) 在 \(\mathrm{Sp}^n(\boldsymbol{a})\) 的一个邻域中等于 1。由构造，它在 \(\mathbf{C}^n - \mathrm{Sp}^n(\boldsymbol{a})\) 上也成立。

引理 4. — 设 \(a \in A^{n}\)。设 \((h, u_{1}, \ldots, u_{n})\) 是一个适配于 a 的序列，\(\omega\) 是相伴微分形式。

\begin{enumerate}
\def\labelenumi{\alph{enumi})}
\tightlist
\item
  对 \(i = 1,2,\ldots,n\)，存在一个微分形式 \(\beta_{i}\)（\(\mathbf{C}^{n}\) 上的、\(n - 1\) 次的、以 A 为系数的微分形式），使得
\end{enumerate}

\[
(z _ {i} - a _ {i}) \omega = d (h \beta_ {i} \wedge d z _ {1} \wedge \dots \wedge d z _ {n});
\]

\begin{enumerate}
\def\labelenumi{\alph{enumi})}
\setcounter{enumi}{1}
\item
  微分形式 \(\omega\) 具有含于 \(h\) 的支集中的紧支集；
\item
  存在一个微分形式 \(\beta\)（\(\mathbf{C}^{n}\) 上的、\(n - 1\) 次的、以 A 为系数的微分形式），使得
\end{enumerate}

\[
(n + 1) h \omega - \omega = d (h \beta \wedge d z _ {1} \wedge \dots \wedge d z _ {n}).
\]

设 \(i\) 是一个满足 \(1 \leqslant i \leqslant n\) 的整数。存在 \(\varepsilon_i \in \{-1, 1\}\)，使得

\[
\varepsilon_ {i} \bigwedge_ {j \neq i} d u _ {j} \wedge d z _ {1} \wedge \dots \wedge d z _ {n} = d z _ {i} \wedge \bigwedge_ {j \neq i} (d u _ {j} \wedge d z _ {j}).
\]

令 \(\beta_{i} = \varepsilon_{i} \bigwedge_{j \neq i} du_{j}\)，于是该公式的左端项是 \(\beta_{i} \wedge dz_{1} \wedge \cdots \wedge dz_{n}\)，且 \(d\beta_{i} = 0\)。这样

\[
d \left(h \beta_ {i} \wedge d z _ {1} \wedge \dots \wedge d z _ {n}\right) = d h \wedge \beta_ {i} \wedge d z _ {1} \wedge \dots \wedge d z _ {n} =
\]

\[
d h \wedge d z _ {i} \wedge \bigwedge_ {j \neq i} (d u _ {j} \wedge d z _ {j}) = (z _ {i} - a _ {i}) \omega ,
\]

这是依公式 (3)，由此得断言 \(a\) )。

由断言 \(a\) ) 及公式 (1)，我们推出关系式

\[
\begin{array}{l} \omega = h \omega + (1 - h) \omega = h \omega + \sum_ {i = 1} ^ {n} (z _ {i} - a _ {i}) u _ {i} \omega \\ \qquad \qquad \qquad = h \omega + \sum_ {i = 1} ^ {n} u _ {i}   d (h \beta_ {i} \wedge d z _ {1} \wedge \dots \wedge d z _ {n}), \end{array}
\]

由此 \(\operatorname{Supp}(\omega) \subset \operatorname{Supp}(h)\)，这就证明了 \(b\) )。

最后，令

\[
\beta = \sum_ {i = 1} ^ {n} \varepsilon_ {i} u _ {i} \bigwedge_ {j \neq i} d u _ {j} = \sum_ {i = 1} ^ {n} u _ {i} \beta_ {i}, \text {   and   } \tau = h \beta d z _ {1} \wedge \dots \wedge d z _ {n}.
\]

我们有 \(d\beta = \sum_{i}du_{i}\wedge \beta_{i}\)，因而

\[
d \beta \wedge d z _ {1} \wedge \dots \wedge d z _ {n} = \sum_ {i} d u _ {i} \wedge d z _ {i} \wedge \bigwedge_ {j \neq i} (d u _ {j} \wedge d z _ {j}) = n \omega .
\]

于是

\[
\begin{array}{l} d \tau = d h \wedge \beta \wedge d z _ {1} \wedge \dots \wedge d z _ {n} + h d \beta \wedge d z _ {1} \wedge \dots \wedge d z _ {n} \\ = \sum_ {i = 1} ^ {n} u _ {i} d h \wedge d z _ {i} \wedge \bigwedge_ {j \neq i} (d u _ {j} \wedge d z _ {j}) + n h \omega \\ = - \sum_ {i = 1} ^ {n} u _ {i} (z _ {i} - a _ {i}) \omega + n h \omega = (h - 1) \omega + n h \omega = (n + 1) h \omega - \omega , \end{array}
\]

给定公式 (3) 与 (1)，由此得 c)。

我们现在来研究微分形式 \(\omega\)（与一个适配于 \(\boldsymbol{a}\) 的序列相伴的）如何作为该序列的函数而变化。我们说两个序列 \((h, u_1, \ldots, u_n)\) 与 \((h', u_1', \ldots, u_n')\)（均适配于 \(\boldsymbol{a}\)）是相联的，如果存在一个微分形式 \(\psi\)（\(n-1\) 次的、\(\mathbf{C}^n\) 上的、以 A 为系数且支包含于 \(h\) 与 \(h'\) 的支集之并中），使得相伴的微分形式 \(\omega\) 与 \(\omega'\) 满足

\[
\omega - \omega^ {\prime} = d (\psi \wedge d z _ {1} \wedge d z _ {2} \wedge \dots \wedge d z _ {n}).
\]

让我们从一个初等变形开始：

引理 5. --- 设 \(\boldsymbol{a} \in \mathrm{A}^{n}\)，设 \((h, u_{1}, \ldots, u_{n})\) 是一个适配于 \(\boldsymbol{a}\) 的序列，并设 \(\omega\) 是相伴微分形式。

设 \(w\) 是一个从 \(\mathbf{C}^n\) 到 \(A\) 中的无穷次可微映射，并设 \(i\) 与 \(j\) 是 1 与 \(n\) 之间的两个不同整数。用下式定义 \(u_1', \ldots, u_n'\)：

\[
\begin{array}{c} u _ {i} ^ {\prime} = u _ {i} + (z _ {j} - a _ {j}) w, \quad u _ {j} ^ {\prime} = u _ {j} - (z _ {i} - a _ {i}) w, \\ u _ {k} ^ {\prime} = u _ {k} \text {for} k \neq i, j. \end{array}
\]

则序列 \((h,u_{1}^{\prime},\ldots,u_{n}^{\prime})\) 适配于 a，且与序列 \((h,u_{1},\ldots,u_{n})\) 相联。

我们记 \(dz = dz_{1} \wedge \cdots \wedge dz_{n}\)。由于

\[
\begin{array}{c} \sum_ {k = 1} ^ {n} (z _ {k} - a _ {k}) u _ {k} ^ {\prime} (\boldsymbol {z}) = \sum_ {k = 1} ^ {n} (z _ {k} - a _ {k}) u _ {k} (\boldsymbol {z}) + w (\boldsymbol {z}) (z _ {j} - a _ {j}) (z _ {i} - a _ {i}) \\ - w (\boldsymbol {z}) (z _ {i} - a _ {i}) (z _ {j} - a _ {j}) = 1 - h (\boldsymbol {z}) \end{array}
\]

对一切 \(z \in C^{n}\)，序列 \((h, u_{1}', \ldots, u_{n}')\) 适配于 a。此外，我们有

\[
\begin{array}{l} d u _ {i} ^ {\prime} \wedge d u _ {j} ^ {\prime} \wedge d z _ {1} \wedge \dots \wedge d z _ {n} = \\ \left(d u _ {i} + w d z _ {j} + (z _ {j} - a _ {j}) d w\right) \wedge \left(d u _ {j} - w d z _ {i} - (z _ {i} - a _ {i}) d w\right) \wedge d \boldsymbol {z} \\ = \left(d u _ {i} \wedge d u _ {j} - (z _ {i} - a _ {i}) d u _ {i} \wedge d w - (z _ {j} - a _ {j}) d u _ {j} \wedge d w\right) \wedge d \boldsymbol {z} \end{array}
\]

因此存在 \(\varepsilon\in\{-1,1\}\)，使得 \(\varepsilon(\omega-\omega')\) 等于

\[
\begin{array}{l} d u _ {i} ^ {\prime} \wedge d u _ {j} ^ {\prime} \wedge \bigwedge_ {k \neq i, j} d u _ {k} ^ {\prime} \wedge d \boldsymbol {z} - d u _ {i} \wedge d u _ {j} \wedge \bigwedge_ {k \neq i, j} d u _ {k} \wedge d \boldsymbol {z} \\ = - \Big ((z _ {i} - a _ {i}) d u _ {i} \wedge d w + (z _ {j} - a _ {j}) d u _ {j} \wedge d w \Big) \wedge \bigwedge_ {k \neq i, j} d u _ {k} \wedge d \boldsymbol {z} \\ = - \Big (\sum_ {k = 1} ^ {n} (z _ {k} - a _ {k})   d u _ {k} \Big) \wedge d w \wedge \bigwedge_ {k \neq i, j} d u _ {k} \wedge d \boldsymbol {z} \end{array}
\]

而把 (2) 考虑在内，它就等于

\[
d h \wedge d w \wedge \bigwedge_ {k \neq i, j} d u _ {k} \wedge d \boldsymbol {z} = d \left(h d w \wedge \bigwedge_ {k \neq i, j} d u _ {k} \wedge d \boldsymbol {z}\right),
\]

由此即得结论。

引理 6. — 设 \(\mathbf{a} \in \mathrm{A}^n\)。所有适配于 \(\mathbf{a}\) 的序列都是相联的。

设 \((h,u_{1},\ldots,u_{n})\) 与 \((h',u_{1}',\ldots,u_{n}')\) 是两个适配于 a 的序列，并用 \(\omega\) 与 \(\omega'\) 表示相伴的微分形式。

我们来定义无穷次可微映射

\[
\begin{array}{c} {w _ {i j} = u _ {i} ^ {\prime} u _ {j} - u _ {i} u _ {j} ^ {\prime},} \\ {s _ {i} = u _ {i} ^ {\prime} h - u _ {i} h ^ {\prime},} \end{array}
\]

\[
\begin{array}{l} 1 \leqslant i \leqslant n, 1 \leqslant j \leqslant n, \\ 1 \leqslant i \leqslant n, \end{array}
\]

于是 \(w_{ji} = -w_{ij}\)，且 \(\operatorname{Supp}(s_i) \subset \operatorname{Supp}(h) \cup \operatorname{Supp}(h')\)。

我们置 \(u_i'' = u_i' - s_i\)，\(\boldsymbol{u} = (u_1, \ldots, u_n)\) 与 \(\boldsymbol{u}'' = (u_1'', \ldots, u_n'')\)。再用 \(\boldsymbol{v}_{ij}\) 表示从 \(\mathbf{C}^n\) 到 \(A^n\) 中的映射，其第 \(i\) 个分量是 \((z_j - a_j) w_{ij}\)，第 \(j\) 个分量是 \((z_i - a_i) w_{ji} = -(z_i - a_i) w_{ij}\)，其余分量均为零。则我们有

\[
\boldsymbol {u} ^ {\prime \prime} = \boldsymbol {u} + \sum_ {i <   j} \boldsymbol {v} _ {i j}.
\]

事实上，对每个满足 \(1 \leqslant k \leqslant n\) 的整数 k，右端的第 k 个分量是

\[
\begin{array}{c} u _ {k} + \sum_ {i = 1} ^ {k - 1} (z _ {i} - a _ {i}) w _ {k i} + \sum_ {j = k + 1} ^ {n} (z _ {j} - a _ {j}) w _ {k j} = u _ {k} + \sum_ {i = 1} ^ {n} (z _ {i} - a _ {i}) w _ {k i} \\ = u _ {k} + u _ {k} ^ {\prime} \sum_ {i = 1} ^ {n} (z _ {i} - a _ {i}) u _ {i} - u _ {k} \sum_ {i = 1} ^ {n} (z _ {i} - a _ {i}) u _ {i} ^ {\prime} \\ = u _ {k} + (1 - h) u _ {k} ^ {\prime} - (1 - h ^ {\prime}) u _ {k} = u _ {k} ^ {\prime} - s _ {k}. \end{array}
\]

由归纳法，我们从引理 5（应用于整数 \(i\) 与 \(j\) 以及映射 \(w_{ij}\)）推出，序列 \((h,u_1'',\ldots,u_n'')\) 适配于 \(\boldsymbol{a}\) 且与 \((h,u_1,\ldots,u_n)\) 相联。设 \(\omega''\) 是与 \((h,u_1'',\ldots,u_n'')\) 相伴的微分形式。由于 \(u_i'' = u_i' - s_i\)，我们有

\[
\begin{array}{r} \omega^ {\prime \prime} - \omega^ {\prime} = d (u _ {1} ^ {\prime} - s _ {1}) \wedge d z _ {1} \wedge \dots \wedge d (u _ {n} ^ {\prime} - s _ {n}) \wedge d z _ {n} - \\ d u _ {1} ^ {\prime} \wedge d z _ {1} \wedge \dots \wedge d u _ {n} ^ {\prime} \wedge d z _ {n}, \end{array}
\]

它可表为如下形式的微分形式的一个线性组合（系数为 1 或 -1）

\[
\xi_ {\mathrm{I} _ {1}, \mathrm{I} _ {2}} = \bigwedge_ {i \in \mathrm{I} _ {1}} d s _ {i} \wedge \bigwedge_ {i \in \mathrm{I} _ {2}} d u _ {i} ^ {\prime} \wedge d z _ {1} \wedge \dots \wedge d z _ {n},
\]

其中 \(I_{1}\)（相应地，\(I_{2}\)）是 \(\{1,\ldots,n\}\) 的一个非空子集（相应地，一个子集）。每个微分形式 \(\xi_{I_{1},I_{2}}\) 也可写成如下形式

\[
d (\widetilde {\psi} \wedge d z _ {1} \wedge \dots \wedge d z _ {n}),
\]

其中微分形式 \(\widetilde{\psi}\) 的支包含于 \(s_i\) 的支集，对一切 \(i \in \mathrm{I}_1\)。由于 \(\mathrm{I}_1\) 非空，这个支集包含于 \(\operatorname{Supp}(h) \cup \operatorname{Supp}(h')\)。因而 \((h, u_1'', \ldots, u_n'')\) 与 \((h', u_1', \ldots, u_n')\) 相联，并且写出如下式子即得引理

\[
\omega - \omega^ {\prime} = (\omega - \omega^ {\prime \prime}) + (\omega^ {\prime \prime} - \omega^ {\prime}).
\]

\subsection{\texorpdfstring{映射 \(\Theta_{a}\) 的构造}{映射 \textbackslash Theta\_\{a\} 的构造}}

设 \(\boldsymbol{a}=(a_{1},\ldots,a_{n})\in\mathrm{A}^{n}\)，并设 U 是 \(\mathrm{Sp}^{n}(\boldsymbol{a})\) 的一个开邻域。设 h 是一个无穷次可微函数，在 \(\mathrm{Sp}^{n}(\boldsymbol{a})\) 的一个邻域中等于 1，且 h 的支集是紧的并包含于 U（I, p.~52, 引理 1）。依 I, p.~54, 引理 3，存在无穷次可微函数 \((u_{1},\ldots,u_{n})\)（\(C^{n}\) 上的、取值于 A 中的），使得序列 \((h,u_{1},\ldots,u_{n})\) 适配于 a。设 \(\omega\) 是相伴的微分形式；它具有包含于 U 中的紧支集（I, p.~54, 引理 4）。存在一个无穷次可微函数 \(\psi\)（在 U 中具有紧支集且取值于 A 中），使得

\[
\omega = \psi d x _ {1} \wedge d y _ {1} \wedge \dots \wedge d x _ {n} \wedge d y _ {n}
\]

其中 \(x_{j} + iy_{j}\) 是 \(C^{n}\)（被等同于 \(R^{2n}\)）上的坐标函数。设 \(\mu\) 是 \(R^{2n}\) 上的 Lebesgue 测度。

设 \(f \in \mathcal{C}(\mathrm{U};\mathrm{A})\)。微分形式 \(f\omega|\mathrm{U}\)（\(\mathrm{U}\) 上的）是连续的且具有紧支集。与这个微分形式相伴的向量测度（VAR, R2, 10.4.3 与 10.4.4）是向量测度 \(f\psi \cdot \mu\)；它是一个以 \(\mu\) 为底的测度（INT, VI, §2, no. 4, def. 4），它的支集是紧的且包含于 \(\omega\) 的支集。这个测度相对于 A 的范数是有界的（INT, VI, §2, no. 4, prop. 8）；我们用 \(\| f\omega\|\) 表示与之相伴的 \(\mathbf{C}^n\) 上的正测度（INT, VI, §2, no. 3, def. 3）。

依 INT, VI, §2, no. 4, prop. 8, b)，我们有

\[
\| f \omega \| = \| f \psi \cdot \mu \| = \| f \psi \| \cdot \mu = \| f \| \| \omega \|.
\]

特别地，微分形式 \(f\omega\) 在 U 上的积分满足

\[
\left\| \int_ {\mathrm{U}} f \omega \right\| \leqslant \int_ {\mathrm{U}} \| f \| \| \omega \|
\]

（INT, VI, §2, no. 3, prop. 5）。

引理 7. --- 对每个函数 \(f \in \mathcal{O}(\mathrm{U};\mathrm{A})\)，积分 \(\int_{\mathrm{U}} f\omega\) 是 A 的一个元素，它只依赖于 \(\boldsymbol{a}\) 以及 \(f\) 在 \(\mathrm{Sp}^{n}(\boldsymbol{a})\) 的一个邻域中的芽。它满足不等式

\[
\left\| \int_ {\mathrm{U}} f \omega \right\| \leqslant \left(\int_ {\mathrm{U}} \| \omega \|\right) \sup _ {z \in \operatorname{Supp} (h)} \| f (z) \|,\tag{4}
\]

并且

\[
\int_ {\mathrm{U}} (a f) \omega = a \int_ {\mathrm{U}} f \omega
\]

对一切 \(a \in \mathbf{A}\)。

我们在上面已经看到，积分对 \(f \in \mathcal{C}(\mathrm{U};\mathrm{A})\) 有定义，并且满足

\[
\left\| \int_ {\mathrm{U}} f \omega \right\| \leqslant \int_ {\mathrm{U}} \| f \| \| \omega \| \leqslant \left(\int_ {\mathrm{U}} \| \omega \|\right) \sup _ {z \in \operatorname{Supp} (h)} \| f (z) \|.
\]

此外，对每个 \(a \in \mathrm{A}\) 与每个 \(f \in \mathcal{C}(\mathrm{U};\mathrm{A})\)，我们有

\[
\int_ {\mathrm{U}} (a f) \omega = a \int_ {\mathrm{U}} f \omega
\]

（INT, VI, §2, no. 2, prop. 2，应用于乘以 a）。

设 \((h', u_{1}', \ldots, u_{n}')\) 是一个适配于 a 的序列，满足 \(\operatorname{Supp}(h') \subset \operatorname{U}\)。设 \(\omega'\) 是相伴的微分形式。依 I, p.~57, 引理 6，存在一个 n-1 次微分形式 \(\psi\)（\(C^{n}\) 上的）（以 A 为系数且支包含于 \(\operatorname{Supp}(h) \cup \operatorname{Supp}(h') \subset \operatorname{U}\) 中），使得

\[
\omega - \omega^ {\prime} = d (\psi \wedge d z _ {1} \wedge \dots \wedge d z _ {n}).
\]

设 \(f \in \mathcal{O}(\mathrm{U};\mathrm{A})\)。由于映射 \(f\) 是全纯的，我们有

\[
d f = \sum_ {i = 1} ^ {n} \frac {\partial f}{\partial z _ {i}} d z _ {i},
\]

（VAR, R2, p.~24, 8.8.9），因而

\[
f \left(\omega - \omega^ {\prime}\right) = f d \left(\psi \wedge d z _ {1} \wedge \dots \wedge d z _ {n}\right) = d \left(f \psi \wedge d z _ {1} \wedge \dots \wedge d z _ {n}\right).
\]

依 Stokes 公式（VAR, R2, p.~48, 11.2.3），我们于是有

\[
\int_ {\mathrm{U}} f (\omega - \omega^ {\prime}) = 0.
\]

于是元素 \(\int_{\mathrm{U}} f\omega\) 不依赖于序列 \((h, u_1, \ldots, u_n)\) 的选取。为完成证明，我们来证 \(\int_{\mathrm{U}} f\omega\) 只依赖于 \(f\) 在 \(\mathrm{Sp}^n(\boldsymbol{a})\) 的一个邻域中的芽。设 U 与 \(\mathrm{U}'\) 是 \(\mathrm{Sp}^n(\boldsymbol{a})\) 的两个开邻域。设 \(f \in \mathcal{O}(\mathrm{U}; \mathrm{A})\) 与 \(f' \in \mathcal{O}(\mathrm{U}', \mathrm{A})\)，使得 \(f\) 与 \(f'\) 在一个开邻域 \(\mathrm{U}''\)（\(\mathrm{Sp}^n(\boldsymbol{a})\) 的）上相等。存在一个映射 \(h\)（从 \(\mathbf{C}^n\) 到 \(\mathbf{C}\) 中的、无穷次可微的），在 \(\mathrm{Sp}^n(\boldsymbol{a})\) 的一个邻域中等于 1，且具有包含于 \(\mathrm{U}''\) 中的紧支集（I, p.~52, 引理 1），并且存在 \((u_1, \ldots, u_n)\)，使得序列 \((h, u_1, \ldots, u_n)\) 适配于 \(\boldsymbol{a}\)（I, p.~54, 引理 3）。设 \(\omega\) 是相伴的微分形式。由于 \(\mathrm{Supp}(\omega) \subset \mathrm{Supp}(h) \subset \mathrm{U}''\)（I, p.~54, 引理 4, b)），我们有

\[
\int_ {\mathrm{U}} f \omega = \int_ {\mathrm {U^ {\prime\prime}}} f \omega = \int_ {\mathrm {U^ {\prime\prime}}} f ^ {\prime} \omega = \int_ {\mathrm {U^ {\prime}}} f ^ {\prime} \omega ,
\]

这就完成了证明。

这个引理表明，存在唯一的 A-线性映射 \(\Theta_{\boldsymbol{a}}\)（从 \(\mathcal{O}(\mathrm{Sp}^{n}(\boldsymbol{a});\mathrm{A})\) 到 A 中的），使得

\[
\Theta_ {\boldsymbol {a}} (f) = \frac {n !}{(2 i \pi) ^ {n}} \int_ {\mathrm{U}} \widetilde {f} \omega\tag{5}
\]

它对每个开集 U 与每个代表元 \(\tilde{f} \in \mathcal{O}(\mathrm{U};\mathrm{A})\)（某个芽 \(f \in \mathcal{O}(\mathrm{Sp}^n (\boldsymbol {a});\mathrm{A})\) 的）成立。线性映射 \(\Theta_{\boldsymbol{a}}\) 依不等式 (4) 以及 EVT, II, p.~29, prop. 5 是连续的。
